\documentclass[lang=cn,12pt]{elegantbook}

\title{6G意图驱动智能决策与网络高精度孪生建模及编排等关键技术研究报告}

\author{叶帅 \& 冷冰}
\institute{中国移动紫金（江苏）创新研究院}
\date{\today}
\version{1.0}

\extrainfo{本文档版权归中国移动紫金（江苏）创新研究院有限公司所有。未经本公司书面许可，任何单位和个人不得复制、摘录、翻译、编辑本文档任何内容。}

\setcounter{tocdepth}{3}

\logo{logo-zjy.png}
\cover{DTN.png}

% 本文档命令
\usepackage{array}
\newcommand{\ccr}[1]{\makecell{{\color{#1}\rule{1cm}{1cm}}}}
\usepackage{float}
\usepackage{tabularx}
\usepackage{xurl}
\usepackage{booktabs}
\usepackage{listings}

\lstset{
    basicstyle=\ttfamily\small,
    breaklines=true,
    columns=fullflexible,
    keepspaces=true,
    frame=single,
    showstringspaces=false
}
% 修改标题页的橙色带
% \definecolor{customcolor}{RGB}{32,178,170}
% \colorlet{coverlinecolor}{customcolor}

\addbibresource{reference.bib}  % 添加您的 .bib 文件

\begin{document}

\maketitle
\frontmatter

\tableofcontents

\mainmatter

\chapter{研究背景与意义}

面向 6G 网络业务多样化、资源异构化、网络自治化和服务化演进趋势，传统基于人工经验和离线仿真的网络管理方式难以满足低时延、高可靠、可解释和低成本试错需求。网络数字孪生通过构建与物理网络实时交互映射的虚拟网络体，为意图理解、智能决策、策略预验证和闭环编排提供可信环境。本文围绕 6G 意图驱动智能决策、网络高精度孪生建模及编排等关键技术，系统分析技术背景、总体架构、关键机制、挑战与演进方向。

\section{6G 网络演进趋势}

随着 5G 网络规模化部署和垂直行业应用不断深化，移动通信网络正从“以连接为中心”的通信基础设施，逐步演进为“连接、感知、计算、智能、安全”深度融合的新型数字基础设施。预计 2030 年及以后，6G 网络将进一步承载沉浸式通信、通感一体化、工业互联网、车联网、低空经济、空天地一体网络、智能制造、智慧城市等多样化业务需求，其网络形态、能力边界和服务模式都将发生深刻变化。


在此过程中，人工智能将不再只是外部辅助工具，而是逐步嵌入无线接入、核心网、承载网、边缘云和终端等多个网络环节，成为网络设计、运行、优化和演进的内生能力。通过 AI 原生化，网络可以实现流量预测、资源调度、干扰协调、故障识别、能耗优化和用户体验保障等智能功能。未来网络将不只是被动承载业务，而是能够根据业务状态和环境变化进行主动感知、主动分析和主动优化。


此外，未来网络能力将以更加模块化、云原生化和服务化的方式组织，通信、感知、计算、数据、智能、安全和编排等能力可以通过标准化接口和 API 向上层应用开放。网络不再只是提供连接服务，而是逐步演进为可调用、可组合、可编排的新型能力平台。面向差异化行业需求，6G 网络需要支持按需定制、灵活编排和动态调整，以适配不同业务在时延、可靠性、带宽、能耗和安全等方面的差异化要求。


这些趋势使未来网络具备更强能力的同时，也对网络感知、建模、决策、验证和编排提出了更高要求，为意图驱动智能决策、网络高精度孪生建模和自动化编排控制等关键技术的发展奠定了需求基础。

\section{传统网管方式局限}

面向 5G 及以前的移动通信网络，网络规划、建设、维护和优化主要依赖人工经验、规则配置、离线仿真和现网验证等方式。在网络规模相对可控、业务类型相对稳定的阶段，这类方法能够满足基本运维需求。但随着 6G 网络向泛在智能、通感算智融合、服务化架构和端到端自治方向演进，传统网络管理方式逐渐难以适应未来网络的复杂性和动态性。


经验驱动难以支撑复杂网络管理。传统网络优化高度依赖专家经验，运维人员通常根据 KPI、告警、用户投诉和路测结果进行问题定位与参数调整。然而，6G 网络中的管理对象将从传统网元、链路和小区扩展到频谱、波束、算力、模型、切片、业务意图和用户体验等多维要素。网络状态空间显著扩大，人工方式难以及时、全面地识别多因素之间的耦合关系，也难以满足实时或准实时闭环优化要求。


参数配置方式难以支撑意图驱动管理。传统网络管理以底层参数配置为主，例如小区发射功率、切换门限、邻区关系、路由策略和切片资源比例等。该方式要求运维人员将业务目标人工转换为具体网络参数。但在 6G 场景下，用户需求往往表现为高层业务意图，例如“保障工业控制低时延高可靠”“提升热点区域用户体验”“在不影响 SLA 的前提下降低能耗”。传统方式缺少从业务意图到网络策略的自动映射机制，难以实现从意图到配置的自动转换。

\section{意图驱动网络价值}

意图驱动网络（Intent-Based Networking, IBN）是面向 6G 自智网络的重要管理模式。其核心思想是将网络管理方式从“人工配置参数”转变为“用户表达目标、系统自动生成策略”。也就是说，用户或业务系统不再直接指定小区功率、切换门限、路由路径、切片资源比例等底层配置，而是提出高层业务目标，例如“保障某区域低时延业务”“提升热点区域用户体验”“降低夜间网络能耗”等，由系统自动完成意图解析、策略生成、孪生验证、编排执行和闭环保障。

意图驱动网络的基本流程可表示为：

Intent→SLA→Policy→Configuration→Pre-verification→Assurance

其中，Intent 表示业务意图，SLA 表示可量化服务指标，Policy 表示网络策略，Configuration 表示具体配置，Verification 表示策略预验证，Assurance 表示持续保障。


6G 网络涉及无线、承载、核心网、边缘云、算力和业务平台等多个网络域，底层参数数量庞大，人工直接配置难度较高。通过意图驱动机制，运维人员只需关注业务目标，系统自动完成目标分解和策略映射，从而降低人工配置复杂度。传统网络优化往往需要人工分析 KPI、告警、用户投诉和业务需求，再手动制定优化方案。意图驱动网络可以结合 AI、知识图谱和规则推理技术，将业务目标快速转换为可执行策略。例如，“保障工业控制业务低时延高可靠”可被转换为时延、可靠性、资源隔离、调度优先级和边缘算力部署等约束条件，并进一步生成网络配置策略。

\section{网络数字孪生支撑}

网络数字孪生（Network Digital Twin, NDT）是支撑6G网络智能化、自治化和服务化演进的重要技术基础。其核心思想是在数字空间中构建与物理网络实时映射、动态交互的虚拟网络体，通过数据采集、模型构建、状态同步、仿真推演和闭环反馈，实现对物理网络状态、行为和性能的分析、预测、验证与优化。与传统网络仿真不同，网络数字孪生不仅能够描述网络当前运行状态，还能够基于实时数据和高精度模型预测未来网络变化，并在虚拟环境中评估不同控制策略可能产生的影响。因此，网络数字孪生可作为6G网络智能决策的重要验证平台，为策略生成、优化选择和风险评估提供可信依据。


网络数字孪生是实现意图驱动闭环优化的重要支撑。意图驱动网络负责理解业务目标，并将其转换为可执行的网络策略，而数字孪生负责在策略作用于物理网络之前，对策略效果进行预验证。其中，业务意图经过解析后形成网络优化目标和候选策略，随后将候选策略加载至数字孪生环境中，通过仿真推演评估其对网络性能、业务体验和资源消耗的影响。只有满足SLA约束和优化目标的策略，才进一步作用于物理网络。具体而言，网络数字孪生在意图驱动闭环中主要承担以下作用：

（1）网络状态镜像：通过实时采集物理网络中的拓扑、性能、业务和资源数据，实现虚实网络状态同步，为智能决策提供准确环境信息；

（2）策略预验证：在数字空间中模拟候选策略执行过程，分析其对时延、可靠性、吞吐量、能耗等指标的影响，降低直接在现网试错带来的风险；

（3）模型动态校准：根据物理网络执行结果持续更新孪生模型，提高数字空间与真实网络之间的一致性；

（4）优化反馈支撑：结合策略执行效果和网络变化趋势，为后续意图解析和策略生成提供反馈，实现持续优化。

因此，网络数字孪生并不是简单复制物理网络，而是构建连接“业务意图—智能决策—策略验证—网络优化”的虚拟闭环环境，为6G网络实现自主感知、自主决策和自主优化提供关键技术支撑。

\chapter{总体架构与技术体系}

面向6G网络业务多样化、资源异构化和网络自治化发展需求，本文构建面向意图驱动智能决策与网络高精度孪生建模的总体技术架构。该架构以网络数字孪生为基础支撑，以业务意图理解和智能决策为核心，以自动化编排和闭环优化为执行手段，实现从业务需求表达、网络状态感知、孪生模型构建、策略生成验证到网络执行反馈的闭环。

本章采用三层架构、三域模型和双闭环机制的总体架构\cite{Sun_DTN_2021}，并在此架构下研究意图驱动智能决策、高精度孪生建模以及智能编排协同机制。

\section{总体架构}

\subsection*{三层架构}

三层为物理网络层、孪生网络层和网络应用层。

物理网络层包括真实运行的网络基础设施，包括无线接入网、核心网、承载网络、边缘计算节点以及终端设备等。该层负责提供真实网络运行环境，并通过数据采集接口向孪生网络层提供网络状态信息，同时接收经过验证后的优化策略。

孪生网络层是数字孪生系统的核心，包括数据域、模型域和管理域。其中：
\begin{itemize}
    \item 数据域负责多源异构网络数据采集、存储、治理和服务；
    \item 模型域负责构建网络拓扑模型、协议模型、业务模型、无线环境模型以及AI预测模型；
    \item 管理域负责孪生实体生命周期管理、模型更新和虚实一致性维护。
\end{itemize}
孪生网络层通过数据驱动和模型驱动方式，实现对物理网络状态的高精度表达。

网络应用层面向业务需求提供智能化网络服务，包括：
\begin{itemize}
    \item 意图解析服务；
    \item 智能决策服务；
    \item 网络优化服务；
    \item 策略验证服务；
    \item 自动化编排服务。
\end{itemize}
该层通过调用孪生网络能力，实现面向业务目标的智能网络管理。

\subsection*{三域架构}

为了实现网络数字孪生的高效运行，孪生网络层进一步划分为数据域、模型域和管理域。

数据域负责采集和管理网络运行过程中的多维数据，包括：
\begin{itemize}
    \item 网络配置数据；
    \item 性能指标数据；
    \item 用户业务数据；
    \item 告警日志数据；
    \item 环境感知数据。
\end{itemize}
通过数据治理和特征提取，形成支撑模型训练和智能分析的数据基础。

模型域负责构建网络数字孪生模型，包括：
\begin{itemize}
    \item 拓扑模型：描述网络节点和连接关系；
    \item 协议模型：描述网络流程和协议行为；
    \item 无线模型：描述信道、干扰和传播特性；
    \item 业务模型：描述业务流量和用户行为；
    \item AI模型：实现状态预测和智能推理。
\end{itemize}

管理域负责孪生实体生命周期管理，包括模型创建、更新、同步、版本管理和状态维护，保证数字孪生网络与物理网络的一致性。

\subsection*{双闭环机制}

面向6G网络自治需求，本文采用由内闭环和外闭环组成的双闭环优化机制。内闭环运行于数字孪生空间内部，通过模型推演、策略搜索和仿真验证，实现优化方案生成。其流程为：
\begin{equation}
State \rightarrow Prediction \rightarrow Decision \rightarrow Simulation \rightarrow Evaluation
\end{equation}


外闭环连接数字孪生网络和物理网络，将经过验证的优化策略下发至真实网络，并根据执行结果更新孪生模型。其流程为：
\begin{equation}
Decision \rightarrow Orchestration \rightarrow Execution \rightarrow Feedback
\end{equation}

\section{意图驱动智能决策技术架构}

意图驱动智能决策面向6G网络业务目标多样化、服务质量差异化和网络资源动态变化等需求，重点解决“如何将高层业务意图自动转换为网络可执行策略”的问题。该架构以业务意图为输入，以SLA约束建模为中间表达，以智能策略生成为核心输出，形成从业务目标理解到网络策略生成的决策链路。

意图驱动智能决策过程可表示为：

\begin{equation}
Intent \rightarrow SLA \rightarrow Policy \rightarrow Candidate\ Strategy
\end{equation}

其中，$Intent$表示业务意图，$SLA$表示可量化服务等级要求，$Policy$表示网络策略规则，$Candidate\ Strategy$表示待验证的候选优化策略。

从功能组成看，意图驱动智能决策架构主要包括意图输入、意图解析、SLA约束建模、策略生成和策略评估五个核心模块。

\begin{enumerate}
    \item \textbf{意图输入模块：}
    负责接收用户、业务系统或上层应用提出的高层业务目标，包括自然语言描述、业务模板和SLA需求等。例如，“保障工业控制业务低时延高可靠”“提升热点区域用户体验”“在满足业务质量前提下降低网络能耗”等。

    \item \textbf{意图解析模块：}
    负责对业务意图进行语义理解和目标分解，识别业务类型、服务目标、约束条件和优先级等关键信息，实现从自然语言或业务需求到结构化意图表示的转换。

    \item \textbf{SLA约束建模模块：}
    负责将业务意图进一步转换为可量化的网络性能指标和约束条件，例如时延、可靠性、吞吐量、丢包率、资源占用和能耗约束等，从而形成可计算、可优化的网络目标。

    \item \textbf{策略生成模块：}
    根据业务目标、SLA约束和当前网络状态，结合规则推理、优化算法、知识图谱和人工智能模型，生成满足业务需求的候选网络策略，包括资源调度、参数优化、切片保障和节能控制等策略。

    \item \textbf{策略评估模块：}
    对生成的候选策略进行初步筛选和优先级排序，判断其是否满足基本约束条件，并输出待进一步验证的候选策略集合，为后续数字孪生环境中的仿真推演和风险评估提供输入。
\end{enumerate}

上述架构重点描述意图驱动智能决策内部的核心处理流程，即“业务意图理解—SLA约束建模—策略生成—候选策略输出”。其中，网络状态数据作为智能决策的重要输入支撑，网络数字孪生则作为候选策略的验证环境，用于在策略作用于物理网络之前完成仿真推演、效果评估和风险分析。具体的意图解析方法、SLA建模方法和智能策略生成机制将在第三章“意图驱动智能决策”中进一步展开。


\section{网络高精度孪生建模与编排机制}

网络高精度孪生建模与编排机制是实现6G网络智能决策和策略验证的重要基础。该机制以物理网络状态精准映射为目标，通过融合多源异构网络数据、知识模型和人工智能模型，构建能够反映网络实体、业务行为、资源状态和性能变化的高精度数字孪生环境。本文提出的网络高精度孪生建模与编排机制具体包括：

\begin{enumerate}
    \item \textbf{高精度孪生建模模块：}
    基于采集的数据构建网络数字孪生模型，实现对网络实体、连接关系、业务行为和资源状态的精准描述。通过融合拓扑模型、无线模型、业务模型、协议模型和AI预测模型，提高数字空间对物理网络的表达能力。

    \item \textbf{数字孪生环境编排模块：}
    面向不同业务需求和验证任务，对孪生环境中的数据资源、模型资源、仿真场景和计算资源进行统一组织与调度，实现针对特定业务场景的快速孪生环境构建。
\end{enumerate}

通过上述机制，网络数字孪生能够形成从“物理网络状态感知—数字空间建模—孪生环境构建—策略仿真验证”的完整流程，为意图驱动智能决策提供高精度、可交互、可验证的虚拟环境。其中，网络高精度孪生建模方法、模型构建技术以及数字孪生环境编排机制将在第四章“网络高精度孪生建模与编排”中进一步展开。

% ============================================================
% 第3章 意图驱动智能决策技术
%
% 建议导言区包含：
% \usepackage{amsmath,amssymb}
% \usepackage{booktabs}
% \usepackage{tabularx}
% \usepackage{array}
% \usepackage{float}
% \usepackage{xcolor}
% \usepackage{listings}
% \usepackage{graphicx}
% \usepackage{enumitem}
% ============================================================

\chapter{意图驱动智能决策技术}
\label{chap:intent_driven_decision}

意图驱动智能决策是实现网络自动化与自智化管控的重要技术环节。其基本思想是以用户、业务系统或运维人员提出的高层业务目标为输入，由网络系统自动完成业务语义理解、网络目标转译、候选策略生成、策略效果评估和策略迭代改进，从而减少传统网络管理中对人工参数配置和运维经验的依赖。与直接面向设备参数的网络控制方式不同，意图驱动决策首先描述用户希望网络达到的业务效果，而不是直接指定应修改哪些无线参数。例如，用户可以提出如下业务需求：

\begin{quote}
“今晚演唱会期间，保障体育场内视频用户的上行体验，
尽量减少能耗，同时不影响语音业务。”
\end{quote}

上述需求同时包含视频业务保障、RAN节能优化和语音业务保护等多个目标，但没有直接给出策略生效时间、目标小区、无线资源配置、调度权重以及业务质量门限等可执行参数。因此，网络系统需要依次回答以下问题：

\begin{enumerate}[label=\arabic*)]
    \item 用户明确表达了哪些业务目标、作用范围和约束条件；
    \item 这些业务语义在RAN侧具体对应哪些网络对象、性能指标和SLA约束；
    \item 为满足上述目标，RAN侧可以生成哪些控制动作和动作参数组合；
    \item 各候选策略能否满足硬约束，其业务收益、资源代价、能耗和风险如何；
    \item 当候选策略不满足要求时，如何根据评估结果调整动作参数或动作组合。
\end{enumerate}

本报告围绕上述问题构建一个面向单小区无线接入网（Radio Access Network, RAN）的意图驱动智能决策原型。
当前原型重点验证从自然语言业务需求到候选RAN策略评估的完整处理链路，不涉及承载网、核心网、用户面功能、边缘计算、多小区协同以及真实设备配置下发。

需要说明的是，本文当前采用的业务规则、SLA门限、动作空间和性能模型系数均为仿真假设。本章实验结果主要用于验证意图驱动决策流程的可运行性、可解释性和相对变化趋势，不代表真实无线网络测量结果。

\section{总体决策流程}

\subsection*{基本处理逻辑}

意图驱动智能决策的核心是完成从业务语言到网络控制策略的逐层转换。本文将该过程划分为业务意图抽取、业务意图转译、候选RAN策略生成以及候选策略评估与反馈优化四个主要环节。其总体处理过程可表示为：

\begin{equation}
\begin{aligned}
\text{Business Requirement}
&\xrightarrow{\mathrm{Intent\ Extraction}}
I_{\mathrm{ext}}
\xrightarrow{\mathrm{Intent\ Translation}}
I_{\mathrm{trans}}^{\mathrm{RAN}}
\\
&\xrightarrow{\mathrm{Policy\ Generation}}
\mathcal{P}^{(0)}
\xrightarrow{\mathrm{Policy\ Evaluation}}
\mathcal{S}^{(0)}
\xrightarrow{\mathrm{Feedback}}
\mathcal{P}^{(1)},
\end{aligned}
\label{eq:intent_decision_pipeline}
\end{equation}

其中：

\begin{itemize}
    \item $I_{\mathrm{ext}}$表示从自然语言需求中得到的结构化业务意图；
    \item $I_{\mathrm{trans}}^{\mathrm{RAN}}$表示完成RAN作用范围、
    网络目标和SLA约束映射后的转译结果；
    \item $\mathcal{P}^{(0)}$表示初始候选RAN策略集合；
    \item $\mathcal{S}^{(0)}$表示候选策略的性能、约束、评分和排序结果；
    \item $\mathcal{P}^{(1)}$表示根据评估反馈生成的下一轮局部改进策略集合。
\end{itemize}

\section{业务意图抽取}

业务意图抽取是意图驱动智能决策链路的起始环节。其目标是从用户、业务系统或运维人员提出的自然语言需求中，识别业务类型、作用范围、目标对象、保障目标、优化方向和保护约束，并将非结构化文本转换为具有统一Schema的结构化业务意图。对于如下业务需求：

\begin{quote}
“今晚演唱会期间，保障体育场内视频用户的上行体验，
尽量减少能耗，同时不影响语音业务。”
\end{quote}

系统需要识别其中包含的时间表达、空间表达、目标用户、目标业务、业务保障目标、节能优化目标和语音保护约束。但是，本阶段不负责确定演唱会的准确起止时间、体育场对应的目标小区、具体SLA门限以及RAN控制参数。上述信息需要在后续业务意图转译和策略生成阶段进一步补全。因此，业务意图抽取主要回答以下两个问题：

\begin{enumerate}[label=\arabic*)]
    \item 用户在原始文本中明确表达了什么；
    \item 哪些信息尚未明确，需要由后续知识库、规则或人工确认进行补充。
\end{enumerate}

业务意图抽取结果可抽象表示为：

\begin{equation}
I_{\mathrm{ext}}
=
\left\{
T,L,O,S,G,C,R,U
\right\},
\label{eq:intent_extraction_representation}
\end{equation}

其中：

\begin{itemize}
    \item $T$表示时间表达；
    \item $L$表示空间表达；
    \item $O$表示意图作用对象；
    \item $S$表示目标业务；
    \item $G$表示业务保障目标和优化目标；
    \item $C$表示业务保护约束；
    \item $R$表示语义关系和逻辑关系；
    \item $U$表示尚未解析或需要补充的信息。
\end{itemize}


\subsection*{问题定义与理论方法}

自然语言业务需求通常具有表达灵活、信息不完整、目标复合和约束隐含等特点。同一业务目标可以有多种表述方式，同一句话中也可能同时包含多个业务目标和保护约束。因此，业务意图抽取并不是简单的关键词识别，而是包含意图分类、实体识别、槽位填充、关系识别和缺失项检测等多个子任务。从理论上看，业务意图抽取可以采用以下几类方法。

\subsubsection*{基于规则和词典的方法}

基于规则的方法通过关键词、正则表达式、领域词典和句式模板，识别自然语言中的业务类型和语义槽位。设输入文本为$x$，规则集合为$\mathcal{R}$，领域关键词集合为$\mathcal{K}$，则规则驱动的抽取过程可表示为：

\begin{equation}
I_{\mathrm{ext}}
=
f_{\mathrm{rule}}
\left(
x;\mathcal{R},\mathcal{K}
\right).
\label{eq:rule_based_extraction_general}
\end{equation}

该类方法具有以下优点：

\begin{itemize}
    \item 处理过程清晰，规则命中结果可解释；
    \item 不依赖大规模标注数据；
    \item 结果稳定，便于单元测试和复现；
    \item 适合业务范围明确、句式相对固定的初步原型。
\end{itemize}

其主要不足是规则覆盖能力有限，对于同义表达、复杂句式和隐含语义的泛化能力较弱。

\subsubsection*{基于监督学习的方法}

对于具有一定规模标注数据的场景，可以采用文本分类和序列标注模型完成意图抽取。多标签意图分类的一般形式为：

\begin{equation}
P(c_k\mid x)
=
\sigma
\left(
\mathbf{w}_k^{T}\mathbf{h}_{\mathrm{CLS}}
+
b_k
\right),
\label{eq:multilabel_intent_classification}
\end{equation}

其中，
$x$表示输入文本，
$c_k$表示第$k$类意图，
$\mathbf{h}_{\mathrm{CLS}}$表示由预训练语言模型得到的文本表示，
$\sigma(\cdot)$表示Sigmoid函数。

当$P(c_k\mid x)$超过相应阈值时，
认为输入文本包含第$k$类意图。
由于一个业务请求可能同时包含业务保障、节能优化和业务保护，
因此应采用多标签分类，而不是单一类别分类。

实体和槽位抽取可采用命名实体识别
（Named Entity Recognition, NER）、
槽位填充（Slot Filling）和BIO序列标注等方法。
例如，可以利用BERT编码器提取上下文表示，
再利用条件随机场
（Conditional Random Field, CRF）
建模相邻标签之间的约束。

需要说明的是，
式\eqref{eq:multilabel_intent_classification}
描述的是学习型意图分类的一般方法，
本文当前原型未训练或调用该类模型。

\subsubsection*{基于大语言模型的方法}

大语言模型可以结合预定义JSON Schema、
Function Calling或Constrained Decoding，
直接输出结构化意图对象。
该方法适合处理长文本、复杂句式和隐含关系，
但需要额外解决以下问题：

\begin{itemize}
    \item 输出字段是否满足既定Schema；
    \item 模型是否生成了原文中不存在的信息；
    \item 置信度是否具有可解释含义；
    \item 网络控制相关高风险字段是否需要人工确认；
    \item 相同输入是否能够获得稳定、一致的输出。
\end{itemize}

因此，在工程系统中，
大语言模型通常需要与规则校验、
领域词典和结构化验证机制结合使用。

\subsubsection*{仿真方法}

本文的重点是验证从自然语言业务需求到RAN候选策略评估的完整闭环，
而不是构建通用自然语言理解模型。
为保证结果可解释、可复现并便于自动化测试，
当前原型采用：

\begin{equation}
\boxed{
\text{有序正则表达式}
+
\text{领域关键词}
+
\text{规则模板}
+
\text{Schema校验}
}
\label{eq:implemented_intent_extraction_method}
\end{equation}

构建规则驱动的业务意图抽取基线。

该方法不调用BERT、CRF、大语言模型、
外部自然语言处理服务或在线API，
整个抽取过程是确定性的，不包含随机采样。


\subsection*{结构化意图表示}

为便于后续业务意图转译，
本文采用统一的数据对象表示业务意图抽取结果。
结构化意图主要包括：

\begin{itemize}
    \item 意图标识和原始文本；
    \item 多标签意图类型；
    \item 字段级语义槽位；
    \item 语义关系和逻辑关系；
    \item 尚未解析的字段；
    \item 抽取方法、信息来源和置信度。
\end{itemize}

结构化意图可表示为：

\begin{equation}
I_{\mathrm{ext}}
=
\left(
ID,
X,
C_I,
E,
R_s,
R_l,
U,
M
\right),
\label{eq:structured_intent_full}
\end{equation}

其中：

\begin{itemize}
    \item $ID$表示意图唯一标识；
    \item $X$表示原始自然语言文本；
    \item $C_I$表示多标签意图类型集合；
    \item $E$表示抽取的语义槽位集合；
    \item $R_s$表示语义关系；
    \item $R_l$表示逻辑关系；
    \item $U$表示待解析项集合；
    \item $M$表示抽取方法、来源和置信度等元数据。
\end{itemize}


\subsection*{多标签意图类型识别}
\label{subsec:implemented_intent_type_detection}

实际业务需求通常同时包含多个目标。
对于本文的演唱会场景，
当前原型支持识别以下三类意图标签：

\begin{equation}
C_I
=
\left\{
c_{\mathrm{assurance}},
c_{\mathrm{energy}},
c_{\mathrm{protection}}
\right\},
\label{eq:implemented_intent_labels}
\end{equation}

分别表示：

\begin{itemize}
    \item \texttt{service\_assurance}：业务保障意图；
    \item \texttt{energy\_optimization}：节能优化意图；
    \item \texttt{service\_protection}：业务保护意图。
\end{itemize}

当前规则中，
当文本同时出现保障语义以及视频或上行语义时，
生成业务保障标签；
当文本出现节能、降低能耗或减少能耗等语义时，
生成节能优化标签；
当文本同时出现语音业务和保护、不得影响等语义时，
生成业务保护标签。

对于示例文本，
可以得到：

\begin{equation}
C_I
=
\left\{
\texttt{service\_assurance},
\texttt{energy\_optimization},
\texttt{service\_protection}
\right\}.
\label{eq:concert_intent_labels}
\end{equation}

多个标签并存说明该业务需求属于复合意图，
其目标不是单纯提升视频业务性能，
而是在满足视频保障和语音保护的前提下，
进一步考虑RAN节能。

如果输入文本不能触发任何受支持的意图标签，
系统不会生成不确定结果，
而是明确返回不支持的意图错误。


\subsection*{语义槽位抽取与规范化}
\label{subsec:semantic_slot_extraction}

意图类型只能说明业务需求属于哪一类，
还不能完整描述业务作用时间、空间、对象和目标。
因此，需要进一步抽取关键语义槽位。

当前原型支持的主要槽位如表
\ref{tab:implemented_intent_slots}所示。

\begin{table}[H]
    \centering
    \caption{当前原型支持的业务意图槽位}
    \label{tab:implemented_intent_slots}
    \small
    \renewcommand{\arraystretch}{1.25}
    \begin{tabularx}{\textwidth}{
        >{\raggedright\arraybackslash}p{0.29\textwidth}
        >{\raggedright\arraybackslash}p{0.28\textwidth}
        >{\raggedright\arraybackslash}X
    }
        \toprule
        槽位类型 & 含义 & 示例 \\
        \midrule

        \texttt{TIME\_EXPRESSION}
        & 时间范围原始表达
        & 今晚演唱会期间
        \\

        \texttt{LOCATION\_EXPRESSION}
        & 空间范围原始表达
        & 体育场内
        \\

        \texttt{TARGET\_OBJECT}
        & 意图作用对象
        & 视频用户
        \\

        \texttt{TARGET\_SERVICE}
        & 目标业务
        & 上行视频业务
        \\

        \texttt{ASSURANCE\_TARGET}
        & 需要保障的目标
        & 上行体验
        \\

        \texttt{OPTIMIZATION\_GOAL}
        & 次级优化目标
        & 降低RAN能耗
        \\

        \texttt{PROTECTED\_SERVICE}
        & 需要保护的业务
        & 语音业务
        \\

        \texttt{MODIFIER}
        & 程度或组合修饰语
        & 尽量、同时
        \\

        \texttt{NEGATION}
        & 非退化或禁止关系
        & 不影响
        \\

        \bottomrule
    \end{tabularx}
\end{table}

示例文本可以标注为：

\begin{quote}
[\textbf{今晚演唱会期间}]$_{\mathrm{TIME}}$，
保障[\textbf{体育场内}]$_{\mathrm{LOCATION}}$
[\textbf{视频用户}]$_{\mathrm{TARGET}}$的
[\textbf{上行体验}]$_{\mathrm{ASSURANCE}}$，
[\textbf{尽量减少能耗}]$_{\mathrm{OPTIMIZATION}}$，
同时[\textbf{不影响语音业务}]$_{\mathrm{CONSTRAINT}}$。
\end{quote}

每个槽位不仅保存规范化结果，
还保留原始文本、信息来源和字段级置信度。
可以表示为：

\begin{equation}
e_j
=
\left(
r_j,
n_j,
s_j,
q_j
\right),
\label{eq:semantic_slot_model}
\end{equation}

其中：

\begin{itemize}
    \item $r_j$表示原始文本片段；
    \item $n_j$表示规范化后的领域值；
    \item $s_j$表示信息来源；
    \item $q_j$表示字段置信度。
\end{itemize}

例如，“视频用户”可以保存为：

\begin{equation}
e_{\mathrm{target}}
=
\left(
\text{视频用户},
\texttt{video\_users},
\texttt{explicit\_text},
q_{\mathrm{target}}
\right).
\end{equation}

原始文本和规范化值必须分开保存。
原始文本用于保留用户直接表达的内容，
规范化值则用于后续程序判断和知识库查询。
这样可以避免将系统推导或术语规范化结果错误地标记为用户原始输入。

当前原型中的规范化值由固定规则和小型领域词典确定。
例如：

\begin{itemize}
    \item “体育场”规范化为\texttt{stadium}；
    \item “视频用户”规范化为\texttt{video\_users}；
    \item “上行视频”规范化为\texttt{uplink\_video}；
    \item “降低能耗”规范化为\texttt{minimize\_ran\_energy}。
\end{itemize}


\subsection*{语义关系与逻辑关系识别}
\label{subsec:semantic_logical_relation_detection}

仅抽取槽位仍不足以表达各字段之间的作用关系。
例如，“语音业务”和“不影响”必须结合起来，
才能表达语音业务是需要保护的对象。

对于示例文本，
当前原型构造以下主要语义关系：

\begin{equation}
\operatorname{Assure}
\left(
\texttt{video\_uplink\_experience}
\right),
\label{eq:assure_relation}
\end{equation}

\begin{equation}
\operatorname{Minimize}
\left(
\texttt{ran\_energy}
\right),
\label{eq:minimize_relation}
\end{equation}

\begin{equation}
\operatorname{Protect}
\left(
\texttt{voice\_service}
\right).
\label{eq:protect_relation}
\end{equation}

时间和地点分别构成意图的时间范围与空间范围：

\begin{equation}
\operatorname{TemporalScope}
\left(
I_{\mathrm{ext}},
T
\right),
\end{equation}

\begin{equation}
\operatorname{SpatialScope}
\left(
I_{\mathrm{ext}},
L
\right).
\end{equation}

除语义关系外，
还需要确定多个目标之间的逻辑关系。

在本文场景中：

\begin{itemize}
    \item 视频业务保障属于主要业务目标；
    \item 降低能耗属于尽可能实现的软优化目标；
    \item 语音业务保护属于强制约束；
    \item “同时”表示多个目标需要共同考虑；
    \item “尽量”表示节能目标不应以破坏硬约束为代价；
    \item “不影响”表示语音质量不能出现不可接受的下降。
\end{itemize}

该复合关系可以表示为：

\begin{equation}
\begin{aligned}
R_l
=
\{
&\operatorname{Primary}
(\texttt{video\_uplink\_assurance}),
\\
&\operatorname{SoftObjective}
(\texttt{ran\_energy\_optimization}),
\\
&\operatorname{HardConstraint}
(\texttt{voice\_service\_protection})
\}.
\end{aligned}
\label{eq:logical_relation_model}
\end{equation}

当前实现通过连接词、否定词、保护关键词和意图标签组合生成上述关系，
并未采用依存句法分析、语义角色标注或大语言模型推理。


\subsection*{规则驱动抽取流程}
\label{subsec:implemented_rule_based_flow}

本文业务意图抽取器的实际处理流程如图式
\eqref{eq:implemented_extraction_flow}所示：

\begin{equation}
\begin{aligned}
x
&\xrightarrow{\mathrm{Normalization}}
x'
\xrightarrow{\mathrm{Regex/Keyword\ Matching}}
M
\\
&\xrightarrow{\mathrm{Label\ Generation}}
C_I
\xrightarrow{\mathrm{Slot\ Construction}}
E
\\
&\xrightarrow{\mathrm{Relation\ Construction}}
R
\xrightarrow{\mathrm{Missing\ Check}}
U
\xrightarrow{\mathrm{Schema\ Validation}}
I_{\mathrm{ext}}.
\end{aligned}
\label{eq:implemented_extraction_flow}
\end{equation}

具体包括以下步骤。

\subsubsection*{文本规范化}

首先去除多余空白，
并统一部分中文标点和文本形式，
减少格式差异对规则匹配的影响。

规范化过程可表示为：

\begin{equation}
x'
=
\operatorname{NormalizeText}(x).
\end{equation}

当前规范化主要服务于规则稳定匹配，
不执行中文分词、词性标注和句法分析。

\subsubsection*{正则和关键词匹配}

系统按照预设顺序匹配：

\begin{itemize}
    \item 时间表达；
    \item 地点表达；
    \item 视频与上行业务表达；
    \item 保障类动词；
    \item 节能或降低能耗表达；
    \item 语音业务表达；
    \item 保护、不得影响等约束表达；
    \item “尽量”“同时”等修饰和组合关系。
\end{itemize}

匹配结果集合记为：

\begin{equation}
M
=
\left\{
m_1,m_2,\ldots,m_K
\right\},
\end{equation}

其中每个$m_k$包含规则类型、匹配文本和文本位置等信息。

\subsubsection*{意图标签生成}

根据匹配结果生成多标签意图集合。

业务保障标签可抽象表示为：

\begin{equation}
c_{\mathrm{assurance}}
=
\mathbb{I}
\left[
m_{\mathrm{assure}}
\land
\left(
m_{\mathrm{video}}
\lor
m_{\mathrm{uplink}}
\right)
\right].
\label{eq:service_assurance_rule}
\end{equation}

节能标签可表示为：

\begin{equation}
c_{\mathrm{energy}}
=
\mathbb{I}
\left[
m_{\mathrm{energy}}
\right].
\label{eq:energy_intent_rule}
\end{equation}

业务保护标签可表示为：

\begin{equation}
c_{\mathrm{protection}}
=
\mathbb{I}
\left[
m_{\mathrm{voice}}
\land
m_{\mathrm{protect}}
\right].
\label{eq:service_protection_rule}
\end{equation}

只要上述标签中至少有一个成立，
系统即可继续构造结构化意图。
如果所有标签均不成立，
则认为该输入不属于当前原型支持的业务范围。

\subsubsection*{槽位和关系构造}

系统根据匹配结果构造时间、地点、业务、目标和约束槽位，
并进一步生成\texttt{assure}、
\texttt{minimize}、
\texttt{protect}、
\texttt{temporal\_scope}和
\texttt{spatial\_scope}等关系。

槽位构造不是开放式语义生成，
而是规则匹配结果到固定Schema字段的确定性映射：

\begin{equation}
E
=
g_{\mathrm{slot}}
\left(
M,\mathcal{D}_{\mathrm{domain}}
\right),
\label{eq:slot_construction}
\end{equation}

其中，
$\mathcal{D}_{\mathrm{domain}}$表示当前原型中的小型领域规范化词典。

\subsubsection*{待解析项生成}

系统根据Schema要求检查当前文本中仍无法确定的字段，
并生成待解析项集合$U$。

待解析项可以表示为：

\begin{equation}
U
=
\left\{
u_j
\mid
\operatorname{Required}(u_j)
\land
\neg\operatorname{Resolved}(u_j)
\right\}.
\label{eq:unresolved_item_generation}
\end{equation}

生成后按照待解析项编码去重，
避免同一问题被多个规则重复加入。


\subsection*{歧义与缺失信息检测}
\label{subsec:intent_missing_information}

自然语言业务需求通常只能表达业务方向，
不能直接提供网络执行所需的全部数据。
因此，系统不能根据不完整文本随意生成小区、门限或资源参数，
而应显式输出待解析项。

对于本文示例，
主要待解析项包括：

\begin{itemize}
    \item \texttt{exact\_time\_period}：
    演唱会准确开始时间和结束时间；
    \item \texttt{ran\_scope}：
    体育场对应的gNB、小区和载波范围；
    \item \texttt{video\_service\_identification}：
    视频用户或视频业务的RAN侧识别规则；
    \item \texttt{uplink\_experience\_slo}：
    上行体验对应的吞吐率、时延和统计口径；
    \item \texttt{voice\_protection\_threshold}：
    语音质量绝对门限和相对基线下降门限；
    \item \texttt{objective\_weights}：
    业务保障、节能、资源和风险之间的具体评分权重。
\end{itemize}

其中，
准确活动时间和RAN作用范围始终由后续转译阶段解析。
即使原文中包含“今晚演唱会期间”和“体育场内”，
意图抽取器也只保存其原始表达，
不会自行计算具体日期或选择目标小区。

待解析项可以按照严重程度划分为：

\begin{itemize}
    \item \textbf{阻断项：}
    若未解析，则无法进入策略生成阶段；
    \item \textbf{非阻断项：}
    可以使用配置默认值，但必须保留来源和假设说明。
\end{itemize}

例如，活动准确时间、目标小区、视频SLO和语音保护门限属于阻断项；
目标权重在当前原型中可以由配置文件提供，因此属于非阻断项。

本文当前实现采用Schema必填项检查和规则判断生成待解析项，
未实现统计置信度拒识、自动澄清对话和人工确认工作流。


\subsection*{置信度与信息来源}
\label{subsec:intent_confidence_source}

为了支持字段追踪，
结构化意图中的每个槽位均保存信息来源和置信度。

信息来源主要用于区分：

\begin{itemize}
    \item 用户原始文本中的显式信息；
    \item 由领域词典得到的规范化结果；
    \item 由规则关系推导得到的结构化语义；
    \item 由后续外部数据源补全的信息。
\end{itemize}

需要特别说明的是，
当前原型中的置信度不是由机器学习模型计算得到的概率，
而是从配置文件中按不同字段类型读取的固定仿真假设值。

因此，当前置信度可以表示为：

\begin{equation}
q_j
=
q_j^{\mathrm{config}},
\label{eq:configured_confidence}
\end{equation}

而不是：

\begin{equation}
q_j
=
P
\left(
e_j\mid x
\right).
\end{equation}

该设计的主要目的是验证以下工程接口：

\begin{itemize}
    \item 字段级置信度能否随结构化结果一同传递；
    \item 后续模块能否保留信息来源；
    \item 是否能够区分显式输入与系统规范化结果。
\end{itemize}

因此，本文实验中的置信度不能解释为自然语言模型的真实预测可信度。


\subsection*{结构化输出}
\label{subsec:intent_structured_output}

经过规则匹配、槽位构造、关系识别和缺失项检查后，
系统输出统一的结构化意图对象。

示例输出如下：

\begin{lstlisting}[language={},caption={业务意图抽取结果示例},
label={lst:structured_intent_example}]
{
  "schemaVersion": "1.0",
  "intentId": "INTENT-CONCERT-UL-001",
  "originalText": "今晚演唱会期间，保障体育场内视频用户的上行体验，尽量减少能耗，同时不影响语音业务。",
  "extractionMethod": "rule_based_v1",

  "intentTypes": [
    {
      "type": "service_assurance"
    },
    {
      "type": "energy_optimization"
    },
    {
      "type": "service_protection"
    }
  ],

  "extractedSemantics": {
    "timeExpression": {
      "rawText": "今晚演唱会期间",
      "normalizedValue": "concert_period",
      "source": "explicit_text",
      "confidence": 0.99
    },

    "locationExpression": {
      "rawText": "体育场内",
      "normalizedValue": "stadium",
      "source": "explicit_text",
      "confidence": 0.99
    },

    "targetObject": {
      "rawText": "视频用户",
      "normalizedValue": "video_users",
      "source": "explicit_text",
      "confidence": 0.97
    },

    "targetService": {
      "rawText": "视频用户",
      "normalizedValue": "uplink_video",
      "source": "normalized_from_explicit_text",
      "confidence": 0.95
    },

    "assuranceTarget": {
      "rawText": "上行体验",
      "normalizedValue": "uplink_experience",
      "source": "explicit_text",
      "confidence": 0.96
    },

    "optimizationGoal": {
      "rawText": "尽量减少能耗",
      "normalizedValue": "minimize_ran_energy",
      "source": "explicit_text",
      "confidence": 0.97
    },

    "protectedService": {
      "rawText": "语音业务",
      "normalizedValue": "voice_service",
      "source": "explicit_text",
      "confidence": 0.96
    }
  },

  "semanticRelations": [
    {
      "relation": "assure",
      "object": "video_uplink_experience"
    },
    {
      "relation": "minimize",
      "object": "ran_energy"
    },
    {
      "relation": "protect",
      "object": "voice_service"
    },
    {
      "relation": "temporal_scope",
      "object": "concert_period"
    },
    {
      "relation": "spatial_scope",
      "object": "stadium"
    }
  ],

  "logicalRelations": {
    "combination": "conjunction",
    "energyOptimizationType": "soft_objective",
    "voiceProtectionType": "hard_constraint"
  },

  "unresolvedItems": [
    {
      "code": "exact_time_period",
      "severity": "blocking"
    },
    {
      "code": "ran_scope",
      "severity": "blocking"
    },
    {
      "code": "video_service_identification",
      "severity": "blocking"
    },
    {
      "code": "uplink_experience_slo",
      "severity": "blocking"
    },
    {
      "code": "voice_protection_threshold",
      "severity": "blocking"
    },
    {
      "code": "objective_weights",
      "severity": "non_blocking"
    }
  ]
}
\end{lstlisting}

上述输出中：

\begin{itemize}
    \item \texttt{originalText}保留用户原始需求；
    \item \texttt{intentTypes}描述复合意图标签；
    \item \texttt{extractedSemantics}保存字段级槽位；
    \item \texttt{semanticRelations}描述动作与对象关系；
    \item \texttt{logicalRelations}描述主目标、软目标和硬约束；
    \item \texttt{unresolvedItems}显式列出待补充信息；
    \item \texttt{extractionMethod}记录当前采用的规则版本。
\end{itemize}

结构化输出需要通过数据模型验证，
包括：

\begin{itemize}
    \item 原始文本不能为空；
    \item 请求时间必须包含时区；
    \item 意图标识不能为空；
    \item 字段类型必须符合Schema；
    \item 待解析项严重程度必须属于支持的枚举；
    \item 输出应能够稳定序列化和反序列化。
\end{itemize}


\subsection*{异常处理与测试验证}
\label{subsec:intent_extraction_validation}

为保证意图抽取结果能够安全进入后续处理链路，
当前原型对非法输入和不支持意图进行显式处理。

主要异常情况包括：

\begin{itemize}
    \item 输入文本为空或仅包含空白字符；
    \item 请求时间缺少时区；
    \item 文本未匹配任何受支持的业务意图；
    \item 输入字段类型或Schema不合法。
\end{itemize}

对于不支持的文本，
系统明确抛出意图不支持错误，
而不是生成空意图或自动猜测业务含义。

当前测试主要覆盖：

\begin{itemize}
    \item 完整演唱会复合意图；
    \item 缺少时间表达；
    \item 缺少地点表达；
    \item 不包含节能目标；
    \item 不包含语音保护目标；
    \item 空文本输入；
    \item 不支持的文本；
    \item 输入时区和字段约束；
    \item 结构化意图序列化；
    \item 意图抽取结果进入端到端流水线。
\end{itemize}

通过上述测试，
可以验证规则抽取结果的确定性、
字段完整性和异常处理一致性。


\subsection*{能力边界}
\label{subsec:intent_extraction_limitations}

本文采用的规则驱动方法适合验证意图驱动决策的完整流程，
但不能等同于通用业务意图理解系统。

其主要局限包括：

\begin{enumerate}[label=\arabic*)]
    \item \textbf{词汇覆盖范围有限。}

    当前关键词和正则主要面向演唱会视频上行保障场景，
    对同义词、口语表达和跨领域业务需求的覆盖能力有限。

    \item \textbf{句式泛化能力有限。}

    当前方法不进行句法分析和上下文建模，
    对长距离依赖、嵌套条件和复杂否定关系的处理能力较弱。

    \item \textbf{不支持上下文推理。}

    当前抽取器不会利用历史对话、
    用户身份、业务工单和网络运行上下文进行联合判断。

    \item \textbf{置信度不是真实预测概率。}

    置信度来自配置文件，
    只能用于验证数据接口，不能用于统计拒识或风险决策。

    \item \textbf{不解析准确时间和网络对象。}

    “今晚演唱会期间”和“体育场内”仅作为原始语义保存，
    准确时间、场馆、小区和用户范围由后续转译阶段确定。

    \item \textbf{不包含自动澄清机制。}

    当前系统能够显式输出待解析项，
    但尚未根据缺失项自动生成追问或触发人工确认。

    \item \textbf{不具备持续学习能力。}

    当前规则不会根据历史运行结果自动更新，
    也没有利用新样本进行在线训练或词典扩展。
\end{enumerate}

未来可以在保留现有Schema和规则校验机制的基础上，
逐步引入领域预训练语言模型、大语言模型、
知识增强抽取和低置信度人工确认机制，
提升复杂业务场景下的泛化能力。

\section{业务意图转译}
\label{sec:business_intent_translation}

业务意图抽取阶段已经从自然语言需求中识别出时间表达、
空间表达、目标对象、目标业务、保障目标、节能目标和语音保护约束，
并形成统一的结构化业务意图。
但是，这些信息仍然停留在业务语义层，
尚不能直接用于RAN候选策略生成。

例如，对于如下业务需求：

\begin{quote}
“今晚演唱会期间，保障体育场内视频用户的上行体验，
尽量减少能耗，同时不影响语音业务。”
\end{quote}

业务意图抽取能够识别“今晚演唱会期间”“体育场内”
“视频用户”“上行体验”“减少能耗”和“保护语音业务”等语义，
但仍不能直接确定：

\begin{itemize}
    \item 演唱会准确的开始时间和结束时间；
    \item 策略提前生效和活动结束后的恢复时间；
    \item 体育场对应的gNB、小区和载波；
    \item 视频用户和语音用户在RAN侧的识别规则；
    \item “上行体验”对应的RAN性能指标；
    \item 视频和语音业务需要满足的量化门限；
    \item 视频保障、节能和语音保护之间的优先级关系。
\end{itemize}

因此，业务意图转译的目标是将结构化业务语义转换为
RAN侧可识别、可计算和可约束的需求描述。
其基本过程可表示为：

\begin{equation}
I_{\mathrm{ext}}
\xrightarrow{\mathrm{Scope\ Resolution}}
S_{\mathrm{RAN}}
\xrightarrow{\mathrm{Objective\ Mapping}}
G_{\mathrm{RAN}}
\xrightarrow{\mathrm{SLA\ Mapping}}
C_{\mathrm{SLA}}
\xrightarrow{\mathrm{Readiness\ Check}}
I_{\mathrm{trans}}^{\mathrm{RAN}},
\label{eq:intent_translation_process}
\end{equation}

其中，
$I_{\mathrm{ext}}$表示业务意图抽取结果，
$S_{\mathrm{RAN}}$表示RAN作用范围，
$G_{\mathrm{RAN}}$表示RAN目标集合，
$C_{\mathrm{SLA}}$表示SLA和无线资源约束，
$I_{\mathrm{trans}}^{\mathrm{RAN}}$表示最终的RAN侧转译结果。


\subsection*{问题定义与理论方法}

业务意图转译处于业务语义层和网络控制层之间，
其核心任务是回答以下问题：

\begin{enumerate}[label=\arabic*)]
    \item 业务意图在何时、在哪个RAN对象上生效；
    \item 业务目标应映射为哪些RAN性能指标；
    \item 各项性能指标需要达到什么门限；
    \item 哪些目标属于主要目标、次级目标和硬约束；
    \item 当前转译结果是否已经具备进入策略生成阶段的条件。
\end{enumerate}

从理论上看，
业务意图转译可以综合采用时间语义解析、
地理实体链接、GIS、网络资源清单、
用户驻留关系查询、RAN知识图谱、SLA模板推理、
历史KPI分析和QoE模型等技术。

一般情况下，业务意图转译可以表示为：

\begin{equation}
I_{\mathrm{trans}}^{\mathrm{RAN}}
=
f_{\mathrm{trans}}
\left(
I_{\mathrm{ext}},
K_{\mathrm{domain}},
D_{\mathrm{external}},
X_{\mathrm{RAN}}
\right),
\label{eq:general_intent_translation}
\end{equation}

其中，
$K_{\mathrm{domain}}$表示RAN领域知识，
$D_{\mathrm{external}}$表示活动、场馆、用户和SLA等外部数据，
$X_{\mathrm{RAN}}$表示当前网络运行状态。

在完整工程系统中，
系统可以根据动态GIS、用户驻留关系、测量报告和网络负载状态，
实时修正目标小区和目标用户范围。
也可以利用知识图谱建立业务目标、RAN指标、控制动作和资源对象之间的映射关系。

本文当前重点是验证意图驱动决策链路的完整性和可复现性，
因此采用：

\begin{equation}
\boxed{
\text{静态事件库}
+
\text{场馆--小区映射库}
+
\text{SLA模板}
+
\text{确定性规则}
}
\label{eq:implemented_translation_method}
\end{equation}

完成面向单小区RAN的业务意图转译。

当前转译模块不调用动态GIS、实时用户数据库、
知识图谱、大语言模型或外部在线服务，
也不包含随机性。


\subsection*{RAN作用范围、目标与SLA映射}

业务意图转译首先确定意图的具体作用范围，
然后将业务目标转换为RAN目标，
最后通过SLA模板为相关指标绑定量化门限。

\subsubsection*{RAN作用范围解析}

当前实现的范围解析链路为：

\begin{equation}
\text{Time/Location Semantics}
\rightarrow
\text{Event Record}
\rightarrow
\text{Venue}
\rightarrow
\text{Single RAN Cell}.
\label{eq:implemented_scope_mapping}
\end{equation}

业务意图抽取只保留“今晚演唱会期间”等原始时间表达，
不直接计算准确日期和时刻。
转译器根据文本中的“演唱会”等活动语义，
选择静态事件类型：

\begin{equation}
EventType
=
\begin{cases}
\texttt{concert},
& \text{文本包含演唱会语义},\\
\varnothing,
& \text{无法识别受支持的活动类型}.
\end{cases}
\label{eq:event_type_detection}
\end{equation}

随后在事件库中执行唯一匹配。
设满足查询条件的事件记录集合为
$\mathcal{E}_{\mathrm{match}}$，
只有当：

\begin{equation}
\left|
\mathcal{E}_{\mathrm{match}}
\right|
=
1
\label{eq:unique_event_match}
\end{equation}

时，系统才能确定准确的活动时间和场馆。

若事件记录匹配数量为0，
说明静态事件库中不存在对应活动；
若匹配数量大于1，
则无法唯一确定目标事件。
在上述情况下，
系统不会任意选择事件记录，
而是保留相应待解析项。

事件记录提供活动时间范围：

\begin{equation}
T_{\mathrm{event}}
=
\left[
T_{\mathrm{start}},
T_{\mathrm{end}}
\right].
\label{eq:event_time_window}
\end{equation}

考虑策略预部署和活动结束后的观察恢复，
本文配置策略前后生效窗口：

\begin{equation}
T_{\mathrm{eff,start}}
=
T_{\mathrm{start}}
-
\Delta T_{\mathrm{pre}},
\label{eq:effective_start_time}
\end{equation}

\begin{equation}
T_{\mathrm{eff,end}}
=
T_{\mathrm{end}}
+
\Delta T_{\mathrm{post}},
\label{eq:effective_end_time}
\end{equation}

其中，
$\Delta T_{\mathrm{pre}}$表示活动开始前的策略预生效时间，
$\Delta T_{\mathrm{post}}$表示活动结束后的观察和恢复时间。
这两个参数均由配置文件提供，
而不是写死在转译代码中。

事件记录还提供场馆标识$VenueId$。
系统根据场馆—小区映射库，
将场馆映射为gNB、小区和载波：

\begin{equation}
VenueId
\rightarrow
gNBId
\rightarrow
CellId
\rightarrow
CarrierId.
\label{eq:venue_to_cell_mapping}
\end{equation}

RAN范围可以表示为：

\begin{equation}
S_{\mathrm{RAN}}
=
\left\{
T_{\mathrm{event}},
T_{\mathrm{effective}},
V,
\mathcal{G},
\mathcal{C},
\mathcal{K},
U_{\mathrm{target}}
\right\},
\label{eq:ran_scope_model}
\end{equation}

其中，
$V$表示场馆标识，
$\mathcal{G}$表示gNB集合，
$\mathcal{C}$表示小区集合，
$\mathcal{K}$表示载波集合，
$U_{\mathrm{target}}$表示目标用户规则。

本文当前严格限定为单小区场景：

\begin{equation}
\left|
\mathcal{C}
\right|
=
1.
\label{eq:single_cell_constraint}
\end{equation}

若场馆不存在对应小区，
或者映射结果包含多个小区，
当前模型将拒绝该结果，
而不会进行动态小区筛选。

需要注意的是，
地点槽位并不是当前静态仓库查询的直接主键。
实际场馆由唯一事件记录确定，
再根据事件记录中的$VenueId$查询目标RAN对象。

对于目标用户范围，
当前实现根据意图槽位生成两类结构化业务规则：

\begin{equation}
Rule_{\mathrm{video}}
=
\texttt{uplink\_video},
\qquad
Rule_{\mathrm{voice}}
=
\texttt{voice}.
\label{eq:target_service_rules}
\end{equation}

目标用户集合可以抽象表示为：

\begin{equation}
U_{\mathrm{target}}
=
\left\{
u
\mid
Cell(u)=C_{\mathrm{target}},
\;
Service(u)
\in
\left\{
\texttt{uplink\_video},
\texttt{voice}
\right\}
\right\}.
\label{eq:target_user_scope}
\end{equation}

当前原型只生成目标用户和业务规则，
并不实际查询UE标识、QoS Flow、5QI、
S-NSSAI、DRB或业务流分类结果。

\subsubsection*{RAN目标与优先级映射}

业务意图中的“保障上行体验”“尽量减少能耗”
和“不影响语音业务”属于业务语言层目标，
需要进一步转换为RAN侧的目标集合：

\begin{equation}
G_{\mathrm{RAN}}
=
\left\{
G_{\mathrm{video}},
G_{\mathrm{energy}},
G_{\mathrm{voice}}
\right\}.
\label{eq:ran_objective_set}
\end{equation}

若结构化意图包含业务保障标签，
则生成视频上行保障目标：

\begin{equation}
G_{\mathrm{video}}
=
\texttt{video\_uplink\_assurance}.
\label{eq:video_assurance_goal}
\end{equation}

该目标主要关注：

\begin{itemize}
    \item 目标视频用户达标比例；
    \item 视频用户低分位上行吞吐率；
    \item 视频上行时延高分位值；
    \item 视频用户平均BLER。
\end{itemize}

视频上行保障被标记为主要目标：

\begin{equation}
Priority
\left(
G_{\mathrm{video}}
\right)
=
\texttt{primary}.
\label{eq:video_goal_priority}
\end{equation}

若结构化意图包含节能优化标签，
则生成RAN节能目标：

\begin{equation}
G_{\mathrm{energy}}
=
\texttt{ran\_energy\_optimization},
\label{eq:ran_energy_goal}
\end{equation}

其优化方向为：

\begin{equation}
\min E_{\mathrm{RAN}}.
\label{eq:minimize_ran_energy}
\end{equation}

节能属于次级目标：

\begin{equation}
Priority
\left(
G_{\mathrm{energy}}
\right)
=
\texttt{secondary}.
\label{eq:energy_goal_priority}
\end{equation}

这意味着系统只能在视频SLA、语音保护和无线资源约束满足的前提下，
尽可能降低RAN能耗。

若结构化意图包含业务保护标签，
则生成语音业务保护目标：

\begin{equation}
G_{\mathrm{voice}}
=
\texttt{voice\_service\_protection}.
\label{eq:voice_protection_goal}
\end{equation}

语音保护被标记为强制约束：

\begin{equation}
Priority
\left(
G_{\mathrm{voice}}
\right)
=
\texttt{mandatory}.
\label{eq:voice_goal_priority}
\end{equation}

因此，当前场景中的目标关系为：

\begin{equation}
\begin{aligned}
P_{\mathrm{priority}}
=
\{
&\operatorname{Primary}
(\texttt{video\_uplink\_assurance}),
\\
&\operatorname{Secondary}
(\texttt{ran\_energy\_optimization}),
\\
&\operatorname{Mandatory}
(\texttt{voice\_service\_protection})
\}.
\end{aligned}
\label{eq:objective_priority_relationship}
\end{equation}

\subsubsection*{SLA约束映射}

RAN目标映射解决的是“关注哪些网络指标”，
SLA约束映射进一步回答“这些指标需要达到什么程度”。

本文当前使用版本化SLA模板，
并通过配置引用读取具体门限。
门限数值既不是由自然语言抽取器生成，
也不是由转译器直接硬编码。

对于目标视频用户集合
$U_{\mathrm{video}}$，
首先定义单用户上行吞吐率门限
$R_{\mathrm{th}}$。
单用户是否达标可表示为：

\begin{equation}
\mathbb{I}_u
=
\mathbb{I}
\left[
R_u^{\mathrm{UL}}
\geq
R_{\mathrm{th}}
\right].
\label{eq:single_video_user_satisfaction}
\end{equation}

目标视频用户达标比例为：

\begin{equation}
\rho_{\mathrm{video}}
=
\frac{1}
{\left|U_{\mathrm{video}}\right|}
\sum_{u\in U_{\mathrm{video}}}
\mathbb{I}_u.
\label{eq:video_satisfaction_ratio}
\end{equation}

当前视频用户达标约束为：

\begin{equation}
\rho_{\mathrm{video}}
\geq
0.95,
\qquad
R_{\mathrm{th}}
=
20~\mathrm{Mbps}.
\label{eq:video_satisfaction_constraint}
\end{equation}

即至少95\%的目标视频用户上行吞吐率达到20 Mbps。

对于视频上行时延，
系统采用用户级P95统计值：

\begin{equation}
D_{\mathrm{video}}^{P95}
=
Q_{0.95}
\left(
\left\{
D_u^{\mathrm{UL}}
\mid
u\in U_{\mathrm{video}}
\right\}
\right),
\label{eq:video_delay_p95}
\end{equation}

并要求：

\begin{equation}
D_{\mathrm{video}}^{P95}
\leq
30~\mathrm{ms}.
\label{eq:video_delay_constraint}
\end{equation}

对于语音业务，
系统同时检查绝对质量和相对质量下降。

语音绝对质量约束为：

\begin{equation}
Q_{\mathrm{voice}}
\geq
0.95.
\label{eq:voice_absolute_constraint}
\end{equation}

定义策略实施后的语音质量下降为：

\begin{equation}
\Delta Q_{\mathrm{voice}}
=
\max
\left(
0,
Q_{\mathrm{voice}}^{0}
-
Q_{\mathrm{voice}}
\right),
\label{eq:voice_quality_degradation}
\end{equation}

其中，
$Q_{\mathrm{voice}}^{0}$表示策略实施前的语音基线质量。

相对下降约束为：

\begin{equation}
\Delta Q_{\mathrm{voice}}
\leq
0.02.
\label{eq:voice_relative_constraint}
\end{equation}

此外，为避免候选策略过度占用无线资源，
策略实施后的PRB利用率需要满足：

\begin{equation}
\rho_{\mathrm{PRB}}
\leq
0.90.
\label{eq:prb_utilization_constraint}
\end{equation}

因此，本文当前形成的硬约束集合为：

\begin{equation}
C_{\mathrm{SLA}}
=
\left\{
\begin{aligned}
&\rho_{\mathrm{video}}\geq0.95,\\
&D_{\mathrm{video}}^{P95}\leq30~\mathrm{ms},\\
&Q_{\mathrm{voice}}\geq0.95,\\
&\Delta Q_{\mathrm{voice}}\leq0.02,\\
&\rho_{\mathrm{PRB}}\leq0.90.
\end{aligned}
\right.
\label{eq:implemented_sla_constraint_set}
\end{equation}

需要强调的是，
上述门限均为本文仿真配置中的研究假设，
不来源于3GPP、ETSI或运营商现网规范，
不能解释为通用业务保障标准或实网部署建议。


\subsection*{本文的静态知识驱动转译方法}

本文转译器的实际处理流程可以表示为：

\begin{equation}
\begin{aligned}
I_{\mathrm{ext}}
&\xrightarrow{\mathrm{Intent\ Label\ Check}}
G_{\mathrm{candidate}}
\xrightarrow{\mathrm{Event\ Lookup}}
T_{\mathrm{event}},V
\\
&\xrightarrow{\mathrm{Venue\ Mapping}}
S_{\mathrm{cell}}
\xrightarrow{\mathrm{User\ Rule}}
U_{\mathrm{target}}
\\
&\xrightarrow{\mathrm{SLA\ Template}}
C_{\mathrm{SLA}}
\xrightarrow{\mathrm{Blocking\ Check}}
I_{\mathrm{trans}}^{\mathrm{RAN}}.
\end{aligned}
\label{eq:implemented_static_translation_flow}
\end{equation}

当前实现主要依赖表
\ref{tab:translation_data_sources}
所示的数据和配置。

\begin{table}[H]
    \centering
    \caption{业务意图转译的数据与配置来源}
    \label{tab:translation_data_sources}
    \small
    \renewcommand{\arraystretch}{1.25}
    \begin{tabularx}{\textwidth}{
        >{\raggedright\arraybackslash}p{0.27\textwidth}
        >{\raggedright\arraybackslash}p{0.30\textwidth}
        >{\raggedright\arraybackslash}X
    }
        \toprule
        转译内容 & 数据或配置来源 & 当前处理方法 \\
        \midrule

        活动类型和准确时间
        & \texttt{event\_database.json}
        & 根据活动语义执行唯一事件匹配
        \\

        策略前后生效窗口
        & \texttt{simulation\_config.yaml}
        & 在事件开始和结束时间基础上增加配置偏移
        \\

        场馆到RAN对象映射
        & \texttt{venue\_cell\_mapping.json}
        & 将场馆静态映射到唯一gNB、小区和载波
        \\

        视频、语音目标用户规则
        & 结构化意图槽位
        & 生成\texttt{uplink\_video}和\texttt{voice}业务规则
        \\

        RAN目标与优先级
        & 结构化意图标签
        & 通过确定性规则生成主要、次级和强制目标
        \\

        SLA约束定义
        & \texttt{sla\_templates.json}
        & 选择演唱会上行保障模板并裁剪所需约束
        \\

        SLA具体门限
        & \texttt{simulation\_config.yaml}
        & 通过\texttt{thresholdRef}延迟解析
        \\

        \bottomrule
    \end{tabularx}
\end{table}

转译器首先根据结构化意图中的多标签，
判断是否存在视频保障、RAN节能和语音保护目标。
不同意图标签决定：

\begin{itemize}
    \item 需要生成哪些RAN目标；
    \item SLA模板中需要保留哪些约束；
    \item 不同目标采用什么优先级。
\end{itemize}

当存在时间槽位时，
系统根据原始文本中的演唱会语义查询事件库，
并要求返回唯一事件记录。
事件记录提供准确活动时间和场馆标识。

随后，
系统利用场馆标识查询场馆—小区映射库，
获得唯一目标小区及相关gNB和载波。
当前模型和静态仓库均要求单小区映射，
多小区结果直接拒绝。

在SLA映射阶段，
当前实现固定选择：

\begin{equation}
Template
=
\texttt{concert\_uplink\_assurance}.
\label{eq:fixed_sla_template}
\end{equation}

随后根据意图内容裁剪模板中的约束：

\begin{itemize}
    \item 存在视频保障时，保留视频用户达标比例和P95时延约束；
    \item 存在语音保护时，保留语音绝对质量和相对下降约束；
    \item PRB利用率上限始终保留。
\end{itemize}

单个SLA约束可以表示为：

\begin{equation}
c_j
=
\left(
m_j,
o_j,
r_j,
s_j
\right),
\label{eq:sla_constraint_model}
\end{equation}

其中，
$m_j$表示指标名称，
$o_j$表示比较运算符，
$r_j$表示门限配置引用，
$s_j$表示数据来源。

例如，视频达标比例约束可表示为：

\begin{equation}
c_{\mathrm{video}}
=
\left(
\texttt{videoUserSatisfactionRatio},
\geq,
\texttt{sla.video.satisfaction\_ratio},
\texttt{sla\_template}
\right).
\end{equation}

采用门限引用而不是在转译器中直接保存数值，
可以避免SLA配置散落在算法代码中，
也便于不同实验之间进行配置版本追踪。

转译过程中，
系统还会更新意图抽取阶段的待解析项。
例如，当事件库和场馆映射成功后，
可以清除：

\begin{itemize}
    \item 准确活动时间；
    \item 场馆及目标小区范围；
    \item SLA门限；
    \item 目标优先级。
\end{itemize}

若仍有无法解决的信息，
则继续保留在转译输出中。


\subsection*{转译输出、准入控制与能力边界}

完成作用范围、目标和约束映射后，
系统输出统一的RAN侧转译结果：

\begin{equation}
I_{\mathrm{trans}}^{\mathrm{RAN}}
=
\left\{
S_{\mathrm{RAN}},
G_{\mathrm{RAN}},
C_{\mathrm{SLA}},
P_{\mathrm{priority}},
U_{\mathrm{unresolved}},
Z_{\mathrm{status}}
\right\}.
\label{eq:translated_ran_intent_model}
\end{equation}

其中，
$Z_{\mathrm{status}}$表示转译状态。

设转译完成后的阻断项集合为：

\begin{equation}
U_{\mathrm{blocking}}
=
\left\{
u
\in
U_{\mathrm{unresolved}}
\mid
Severity(u)=\texttt{blocking}
\right\}.
\label{eq:blocking_unresolved_set}
\end{equation}

转译状态定义为：

\begin{equation}
Z_{\mathrm{status}}
=
\begin{cases}
\texttt{ready},
&
U_{\mathrm{blocking}}
=
\varnothing,
\\[2mm]
\texttt{blocked},
&
U_{\mathrm{blocking}}
\neq
\varnothing.
\end{cases}
\label{eq:translation_status}
\end{equation}

只有当：

\begin{equation}
Z_{\mathrm{status}}
=
\texttt{ready}
\label{eq:translation_ready_condition}
\end{equation}

时，转译结果才允许进入候选策略生成模块。

若状态为\texttt{blocked}，
系统不会使用默认值掩盖关键字段缺失，
而是终止当前流水线，
并返回明确的阻断原因。

一个简化的转译输出示例如下：

\begin{lstlisting}[language={},
caption={业务意图转译结果示例},
label={lst:translated_ran_intent_example}]
{
  "schemaVersion": "1.0",
  "intentId": "INTENT-CONCERT-UL-001",

  "scope": {
    "eventTime": {
      "startTime": "2026-07-29T19:30:00+08:00",
      "endTime": "2026-07-29T22:30:00+08:00",
      "source": "event_database"
    },

    "effectiveTime": {
      "startTime": "2026-07-29T19:15:00+08:00",
      "endTime": "2026-07-29T22:45:00+08:00",
      "source": "event_database+simulation_config"
    },

    "venueId": "STADIUM-001",

    "ranObjects": {
      "gnbIds": [
        "GNB-001"
      ],
      "cellId": "CELL-101",
      "carrierIds": [
        "NR-CARRIER-01"
      ],
      "source": "venue_cell_mapping"
    },

    "targetUsers": [
      {
        "serviceType": "uplink_video",
        "scope": "target_cell"
      },
      {
        "serviceType": "voice",
        "scope": "target_cell"
      }
    ]
  },

  "ranObjectives": [
    {
      "name": "video_uplink_assurance",
      "priority": "primary"
    },
    {
      "name": "ran_energy_optimization",
      "priority": "secondary"
    },
    {
      "name": "voice_service_protection",
      "priority": "mandatory"
    }
  ],

  "slaConstraints": [
    {
      "metric": "videoUserSatisfactionRatio",
      "operator": ">=",
      "thresholdRef": "sla.video.satisfaction_ratio",
      "source": "sla_template:concert_uplink_assurance"
    },
    {
      "metric": "videoDelayP95Ms",
      "operator": "<=",
      "thresholdRef": "sla.video.delay_p95_ms",
      "source": "sla_template:concert_uplink_assurance"
    },
    {
      "metric": "voiceQualityScore",
      "operator": ">=",
      "thresholdRef": "sla.voice.quality_min",
      "source": "sla_template:concert_uplink_assurance"
    },
    {
      "metric": "voiceQualityDegradation",
      "operator": "<=",
      "thresholdRef": "sla.voice.maximum_degradation",
      "source": "sla_template:concert_uplink_assurance"
    },
    {
      "metric": "prbUtilization",
      "operator": "<=",
      "thresholdRef": "resources.maximum_prb_utilization",
      "source": "sla_template:concert_uplink_assurance"
    }
  ],

  "objectiveRelationship": {
    "primaryObjective": "video_uplink_assurance",
    "secondaryObjective": "ran_energy_optimization",
    "mandatoryConstraints": [
      "video_sla_satisfied",
      "voice_sla_not_degraded",
      "prb_limit_satisfied"
    ]
  },

  "unresolvedItems": [],
  "translationStatus": "ready"
}
\end{lstlisting}

上述输出中，
活动时间来源于事件库，
生效窗口来源于事件时间和配置偏移，
目标小区和载波来源于场馆映射库，
业务目标来源于结构化意图，
SLA约束来源于版本化模板，
具体门限通过配置引用解析。

当前静态知识驱动方法的主要优点是：

\begin{itemize}
    \item 转译过程确定，便于复现和测试；
    \item 每个字段均可记录来源；
    \item SLA门限与算法代码解耦；
    \item 阻断项和准入条件清晰；
    \item 输出可以直接作为后续策略生成模块的输入。
\end{itemize}

但该方法也存在以下局限：

\begin{enumerate}[label=\arabic*)]
    \item 只支持静态演唱会场景和固定事件类型；
    \item 不支持通用自然语言相对时间解析；
    \item 地点槽位不直接作为GIS或地理实体查询键；
    \item 严格限定单小区，不支持多小区和邻区协同；
    \item 不使用用户测量报告、驻留关系和实时负载修正范围；
    \item 只生成业务规则，不执行真实UE或QoS Flow识别；
    \item 固定使用演唱会上行保障SLA模板；
    \item 不使用知识图谱、学习模型或动态SLA推理；
    \item 不读取实时RAN状态。
\end{enumerate}

需要特别说明的是，
当前转译阶段只负责形成RAN作用范围、目标和约束，
实时RAN状态在后续候选策略性能评估阶段进入处理链路。

经过业务意图转译，
系统完成了：

\begin{equation}
\text{结构化业务语义}
\longrightarrow
\text{RAN作用范围}
\longrightarrow
\text{RAN目标}
\longrightarrow
\text{SLA与资源约束}.
\label{eq:intent_to_ran_requirement}
\end{equation}

此时，系统已经明确了“何时、在哪个小区、
针对哪些用户、需要满足哪些目标和约束”，
但尚未确定具体采用哪些RAN动作及其参数。

\section{网络策略生成}
\label{sec:network_policy_generation}

业务意图转译完成后，
系统已经获得策略生效时间、目标小区、目标用户规则、
RAN业务目标、SLA约束和目标优先级。
但是，这些信息仍然只描述了网络需要达到的目标，
尚未明确RAN侧应执行哪些控制动作以及动作参数如何取值。

网络策略生成的任务是将RAN侧目标和约束转换为
多个结构化候选策略，
为后续性能评估、硬约束检查和策略排序提供输入。

对于演唱会视频上行保障场景，
业务意图转译结果已经明确：

\begin{itemize}
    \item 目标小区和策略生效时间；
    \item 目标视频用户和语音业务规则；
    \item 视频上行保障为主要目标；
    \item RAN节能为次级优化目标；
    \item 语音业务保护和PRB利用率上限为硬约束。
\end{itemize}

网络策略生成需要进一步确定：

\begin{itemize}
    \item 采用上行PRB预留还是调度权重调整；
    \item 是否同时使用多个业务保障动作；
    \item 是否加入语音业务资源保护；
    \item 是否加入RAN节能动作；
    \item 不同动作参数采用哪些离散取值；
    \item 哪些动作组合违反能力、冲突或静态资源约束。
\end{itemize}

本文当前的初始策略生成过程可以表示为：

\begin{equation}
I_{\mathrm{trans}}^{\mathrm{RAN}}
+
\mathcal{T}
+
\boldsymbol{\Theta}
+
\mathcal{C}_{\mathrm{static}}
\longrightarrow
\mathcal{P}^{(0)},
\label{eq:initial_policy_generation_process}
\end{equation}

其中：

\begin{itemize}
    \item $I_{\mathrm{trans}}^{\mathrm{RAN}}$表示业务意图转译结果；
    \item $\mathcal{T}$表示策略模板集合；
    \item $\boldsymbol{\Theta}$表示动作参数离散空间；
    \item $\mathcal{C}_{\mathrm{static}}$表示静态能力和冲突约束；
    \item $\mathcal{P}^{(0)}$表示第0轮初始候选策略集合。
\end{itemize}

需要说明的是，
当前代码中的初始策略生成器仅接收业务意图转译结果，
不直接接收或读取实时RAN状态。
当前网络负载、用户吞吐率、时延、语音质量和能耗等状态，
将在下一节的候选策略性能评估阶段进入处理链路。


\subsection*{问题定义与理论方法}

从理论上看，
网络策略生成可以采用规则规划、案例推理、
组合搜索、数学规划、贝叶斯优化、
强化学习和大语言模型辅助生成等方法。

一般情况下，
候选策略生成可以抽象表示为：

\begin{equation}
\mathcal{P}_{\mathrm{candidate}}
=
f_{\mathrm{policy}}
\left(
I_{\mathrm{trans}}^{\mathrm{RAN}},
X_{\mathrm{RAN}},
\mathcal{A}_{\mathrm{RAN}},
K_{\mathrm{policy}}
\right),
\label{eq:general_policy_generation}
\end{equation}

其中：

\begin{itemize}
    \item $X_{\mathrm{RAN}}$表示当前RAN状态；
    \item $\mathcal{A}_{\mathrm{RAN}}$表示可用RAN动作集合；
    \item $K_{\mathrm{policy}}$表示目标、动作和约束之间的策略知识。
\end{itemize}

在更复杂的系统中，
可以根据当前负载、性能差距和历史经验动态选择动作，
也可以通过优化求解器或强化学习搜索连续动作空间。

本文当前重点是验证从RAN目标到候选策略评估的完整链路，
并保证候选生成过程可解释、确定和可复现。
因此采用：

\begin{equation}
\boxed{
\text{规则映射}
+
\text{策略模板}
+
\text{离散参数笛卡尔积}
+
\text{静态约束过滤}
}
\label{eq:implemented_policy_generation_method}
\end{equation}

生成候选RAN策略。

当前生成过程不使用：

\begin{itemize}
    \item 遗传算法；
    \item 贝叶斯优化；
    \item 强化学习；
    \item 整数规划或非线性规划；
    \item 大语言模型；
    \item 随机采样。
\end{itemize}


\subsection*{RAN动作原语与策略模板}

候选策略由一个或多个RAN动作原语组成。
动作原语用于描述控制方向、控制对象和动作参数，
但不直接对应特定厂商设备命令。

单个动作可表示为：

\begin{equation}
a_j
=
\left(
Type_j,
Target_j,
Parameter_j
\right),
\label{eq:ran_action_model}
\end{equation}

其中：

\begin{itemize}
    \item $Type_j$表示动作类型；
    \item $Target_j$表示目标小区、视频用户或语音用户；
    \item $Parameter_j$表示具体动作参数。
\end{itemize}

当前原型实际实现表
\ref{tab:implemented_ran_actions}
所示的四类RAN动作。

\begin{table}[H]
    \centering
    \caption{RAN动作原语}
    \label{tab:implemented_ran_actions}
    \small
    \renewcommand{\arraystretch}{1.25}
    \begin{tabularx}{\textwidth}{
        >{\raggedright\arraybackslash}p{0.32\textwidth}
        >{\raggedright\arraybackslash}p{0.27\textwidth}
        >{\raggedright\arraybackslash}X
    }
        \toprule
        动作类型 & 参数 & 主要作用 \\
        \midrule

        \texttt{UL\_RESOURCE\_RESERVE}
        & \texttt{reservedPrbRatio}
        & 为目标视频用户预留一定比例的上行PRB
        \\

        \texttt{SCHEDULING\_WEIGHT\_ADJUST}
        & \texttt{videoSchedulingWeight}
        & 提高目标视频业务的调度权重
        \\

        \texttt{VOICE\_RESOURCE\_PROTECT}
        & \texttt{minimumVoicePrbRatio}
        & 为语音业务设置最低资源保护比例
        \\

        \texttt{RAN\_ENERGY\_SAVING}
        & \texttt{energySavingIntensity}
        & 对RAN节能控制强度进行抽象建模
        \\

        \bottomrule
    \end{tabularx}
\end{table}

当前代码未实现上行功率控制、负载均衡、
载波休眠、多小区协调和切换参数调整等动作。
这些动作可以作为后续扩展，
但不属于本文当前仿真实现。

在动作原语基础上，
系统根据业务意图转译结果选择策略模板。
策略模板描述一类候选策略包含哪些动作，
但不直接确定动作参数。

当前模板主要包括：

\begin{itemize}
    \item 基线模板；
    \item 调度保障模板；
    \item 资源保障模板；
    \item 联合保障模板；
    \item 语音保护模板；
    \item 带节能动作的保障模板；
    \item 独立节能优化模板。
\end{itemize}

基线模板不包含主动控制动作，
用于表示不增加业务保障和节能控制的参考策略：

\begin{equation}
A_{\mathrm{baseline}}
=
\varnothing.
\label{eq:baseline_action_set}
\end{equation}

当存在视频上行保障目标时，
系统生成调度保障模板：

\begin{equation}
A_{\mathrm{scheduling}}
=
\left\{
a_{\mathrm{schedule}}
\right\}.
\label{eq:scheduling_template}
\end{equation}

若同时存在语音保护目标，
则模板中加入语音保护动作：

\begin{equation}
A_{\mathrm{scheduling+voice}}
=
\left\{
a_{\mathrm{schedule}},
a_{\mathrm{voice}}
\right\}.
\label{eq:scheduling_voice_template}
\end{equation}

资源保障模板表示：

\begin{equation}
A_{\mathrm{resource}}
=
\left\{
a_{\mathrm{resource}}
\right\},
\label{eq:resource_template}
\end{equation}

若需要保护语音业务，
则表示为：

\begin{equation}
A_{\mathrm{resource+voice}}
=
\left\{
a_{\mathrm{resource}},
a_{\mathrm{voice}}
\right\}.
\label{eq:resource_voice_template}
\end{equation}

联合保障模板同时包含PRB资源预留和视频调度权重调整：

\begin{equation}
A_{\mathrm{joint}}
=
\left\{
a_{\mathrm{resource}},
a_{\mathrm{schedule}}
\right\}.
\label{eq:joint_assurance_template}
\end{equation}

存在语音保护要求时，
联合保障模板扩展为：

\begin{equation}
A_{\mathrm{joint+voice}}
=
\left\{
a_{\mathrm{resource}},
a_{\mathrm{schedule}},
a_{\mathrm{voice}}
\right\}.
\label{eq:joint_voice_template}
\end{equation}

当意图中包含RAN节能目标时，
系统还会生成带节能动作的扩展模板：

\begin{equation}
A_{\mathrm{assurance+energy}}
=
A_{\mathrm{assurance}}
\cup
\left\{
a_{\mathrm{energy}}
\right\}.
\label{eq:assurance_energy_template}
\end{equation}

因此，策略模板实现了：

\begin{equation}
\text{RAN目标}
\longrightarrow
\text{动作组合结构}.
\label{eq:goal_to_template_mapping}
\end{equation}

该过程只决定动作组合，
不判断策略实施后是否满足SLA。


\subsection*{离散参数组合与静态过滤}

策略模板确定动作结构后，
还需要为每个动作赋予具体参数。

当前系统从动作库读取动作定义，
再根据动作参数中的
\texttt{valueSetRef}
查询配置文件中的离散取值集合。

设：

\begin{itemize}
    \item $\Theta_r$表示PRB预留比例集合；
    \item $\Theta_w$表示视频调度权重集合；
    \item $\Theta_v$表示语音最低资源比例集合；
    \item $\Theta_e$表示节能强度集合。
\end{itemize}

则完整离散参数空间可以表示为：

\begin{equation}
\boldsymbol{\Theta}
=
\Theta_r
\times
\Theta_w
\times
\Theta_v
\times
\Theta_e.
\label{eq:full_discrete_action_space}
\end{equation}

对于某个只包含部分动作的模板，
系统仅对该模板实际包含的动作参数执行组合。

若模板$T_k$包含$J_k$个动作，
每个动作的参数值集合分别为
$\Theta_{k,1},\Theta_{k,2},\ldots,\Theta_{k,J_k}$，
则候选参数组合为：

\begin{equation}
\Omega_k
=
\Theta_{k,1}
\times
\Theta_{k,2}
\times
\cdots
\times
\Theta_{k,J_k}.
\label{eq:template_cartesian_product}
\end{equation}

当前代码通过笛卡尔积穷举模板内的全部离散参数组合：

\begin{equation}
\Omega_k
=
\operatorname{CartesianProduct}
\left(
\Theta_{k,1},
\Theta_{k,2},
\ldots,
\Theta_{k,J_k}
\right).
\label{eq:cartesian_product_generation}
\end{equation}

例如，
当联合保障模板包含PRB预留、调度权重和语音保护时，
其参数组合可表示为：

\begin{equation}
\Omega_{\mathrm{joint}}
=
\Theta_r
\times
\Theta_w
\times
\Theta_v.
\label{eq:joint_parameter_space}
\end{equation}

加入节能动作后，
参数空间扩展为：

\begin{equation}
\Omega_{\mathrm{joint+energy}}
=
\Theta_r
\times
\Theta_w
\times
\Theta_v
\times
\Theta_e.
\label{eq:joint_energy_parameter_space}
\end{equation}

笛卡尔积能够完整覆盖当前离散动作空间，
但候选数量会随动作数和参数档位快速增加。
因此，系统在生成过程中执行静态约束过滤、
内容去重和候选数量限制。

静态约束过滤主要包括：

\begin{enumerate}[label=\arabic*)]
    \item 动作类型不能重复；
    \item 动作必须存在于当前动作库；
    \item 动作参数Schema必须完整匹配；
    \item 参数值必须属于配置允许的离散空间；
    \item 动作目标必须与当前单小区范围一致；
    \item 动作之间不能存在明确冲突；
    \item 资源预留比例不能明显超过静态资源上限。
\end{enumerate}

对于PRB资源预留和语音资源保护，
当前采用以下静态资源约束：

\begin{equation}
r+v
\leq
\rho_{\mathrm{PRB}}^{\max},
\label{eq:static_prb_resource_constraint}
\end{equation}

其中：

\begin{itemize}
    \item $r$表示视频业务PRB预留比例；
    \item $v$表示语音业务最低PRB比例；
    \item $\rho_{\mathrm{PRB}}^{\max}$表示配置中的最大PRB利用率。
\end{itemize}

需要区分的是，
式\eqref{eq:static_prb_resource_constraint}
只是生成阶段的静态资源检查，
并不代表策略实施后的实际PRB利用率一定满足约束。
策略后的实际PRB利用率需要由性能评估模型结合当前RAN状态计算。

当前实现还定义了部分语义冲突规则。
例如：

\begin{itemize}
    \item 存在语音保护目标时，
    若采用较高视频调度权重却不包含语音保护动作，
    则该组合被视为冲突；
    \item 最大PRB预留比例与最大节能强度同时出现时，
    该动作组合被视为不协调。
\end{itemize}

这些规则属于当前仿真假设，
主要用于验证策略冲突检查机制。

通过静态检查的策略还需要进行内容去重。
设候选策略$P_i$的动作集合为$A_i$，
系统对每个动作的类型、目标和参数进行规范化排序，
构造内容签名：

\begin{equation}
Signature(P_i)
=
\operatorname{Canonical}
\left(
A_i
\right).
\label{eq:policy_deduplication_signature}
\end{equation}

若两个模板产生完全相同的动作类型、目标和参数，
则：

\begin{equation}
Signature(P_i)
=
Signature(P_j),
\label{eq:duplicate_policy_condition}
\end{equation}

系统只保留稳定排序中首次出现的策略。

去重键不包含模板名称和策略标识，
因此不同模板生成相同策略内容时不会产生重复候选。

候选策略按照模板优先级、模板标识和参数顺序稳定排序。
策略标识按最终保留顺序生成：

\begin{equation}
PolicyId_i
=
\texttt{RAN-POLICY-R}
\Vert n
\Vert
\texttt{-}
\Vert i,
\label{eq:policy_id_generation}
\end{equation}

其中，
$n$表示生成轮次，
$i$表示当前轮中的稳定序号。

对于初始策略生成，
$n=0$。

为控制单轮候选数量，
系统设置：

\begin{equation}
\left|
\mathcal{P}^{(n)}
\right|
\leq
N_{\max},
\label{eq:maximum_candidates_constraint}
\end{equation}

其中，
$N_{\max}$由配置项
\texttt{maximum\_candidates\_per\_round}
给出。

当候选数超过上限时，
当前实现按照稳定顺序保留前$N_{\max}$个策略：

\begin{equation}
\mathcal{P}_{\mathrm{retained}}
=
\operatorname{Prefix}
\left(
\mathcal{P}_{\mathrm{valid}},
N_{\max}
\right).
\label{eq:stable_prefix_truncation}
\end{equation}

当前实现不是随机抽样，
也没有根据性能差距、参数覆盖性或历史成功率重新选择候选。


\subsection*{策略输出与能力边界}

经过模板选择、离散参数组合、静态过滤、
内容去重和数量限制后，
系统输出候选RAN策略集合：

\begin{equation}
\mathcal{P}^{(0)}
=
\left\{
P_1^{(0)},
P_2^{(0)},
\ldots,
P_M^{(0)}
\right\}.
\label{eq:initial_candidate_policy_set}
\end{equation}

单个候选策略可以表示为：

\begin{equation}
P_i
=
\left\{
ID_i,
Scope_i,
A_i,
\Theta_i,
G_i,
C_i,
Trigger_i,
Rollback_i,
Meta_i
\right\},
\label{eq:candidate_policy_model}
\end{equation}

其中：

\begin{itemize}
    \item $ID_i$表示策略标识；
    \item $Scope_i$表示生效时间、目标小区和目标用户规则；
    \item $A_i$表示动作集合；
    \item $\Theta_i$表示动作参数；
    \item $G_i$表示策略对应的RAN目标；
    \item $C_i$表示策略需要满足的硬约束；
    \item $Trigger_i$表示策略触发条件；
    \item $Rollback_i$表示回退条件；
    \item $Meta_i$表示生成轮次、父策略和生成原因等元数据。
\end{itemize}

一个简化的候选策略输出示例如下：

\begin{lstlisting}[language={},
caption={候选RAN策略示例},
label={lst:candidate_policy_example}]
{
  "policyId": "RAN-POLICY-R0-0016",
  "intentId": "INTENT-CONCERT-UL-001",
  "generationRound": 0,
  "parentPolicyId": null,

  "scope": {
    "cellId": "CELL-101",
    "startTime": "2026-07-29T19:15:00+08:00",
    "endTime": "2026-07-29T22:45:00+08:00",
    "targetUserRules": [
      "uplink_video",
      "voice"
    ]
  },

  "objectives": [
    "video_uplink_assurance",
    "ran_energy_optimization",
    "voice_service_protection"
  ],

  "actions": [
    {
      "type": "UL_RESOURCE_RESERVE",
      "target": "target_video_users",
      "parameters": {
        "reservedPrbRatio": 0.20
      }
    },
    {
      "type": "SCHEDULING_WEIGHT_ADJUST",
      "target": "target_video_users",
      "parameters": {
        "videoSchedulingWeight": 1.50
      }
    },
    {
      "type": "VOICE_RESOURCE_PROTECT",
      "target": "voice_users",
      "parameters": {
        "minimumVoicePrbRatio": 0.15
      }
    },
    {
      "type": "RAN_ENERGY_SAVING",
      "target": "target_cell",
      "parameters": {
        "energySavingIntensity": 0.10
      }
    }
  ],

  "hardConstraints": [
    "videoUserSatisfactionRatio >= 0.95",
    "videoDelayP95Ms <= 30",
    "voiceQualityScore >= 0.95",
    "voiceQualityDegradation <= 0.02",
    "prbUtilization <= 0.90"
  ],

  "triggerConditions": [
    "effective_time_reached"
  ],

  "rollbackConditions": [
    "voice_sla_violation",
    "prb_overload",
    "effective_time_ended"
  ],

  "generationReason": "initial_template_grid",
  "status": "pending_evaluation"
}
\end{lstlisting}

需要说明的是，
上述数值只是结构示例，
实际参数值由当前配置中的离散动作空间决定。

当前策略生成方法具有以下优点：

\begin{itemize}
    \item 模板和动作含义清晰，生成结果可解释；
    \item 参数全部来自版本化配置，便于实验复现；
    \item 笛卡尔积能够覆盖当前离散动作空间；
    \item 静态过滤能够提前删除明显非法组合；
    \item 内容签名保证候选去重和稳定标识；
    \item 整个过程确定，不受随机种子影响。
\end{itemize}

但当前实现也存在以下能力边界：

\begin{enumerate}[label=\arabic*)]
    \item \textbf{初始生成不使用实时RAN状态。}

    当前生成器不根据网络负载、
    当前吞吐率和时延进行模板选择或参数剪枝。
    同一转译结果在不同网络状态下生成相同的初始候选集合。

    \item \textbf{动作类型有限。}

    当前只实现PRB预留、调度权重、
    语音资源保护和抽象节能四类动作，
    不包含功率控制、负载均衡、载波休眠和多小区协同。

    \item \textbf{采用离散穷举。}

    当前只搜索预定义参数档位，
    不能优化连续参数，
    也不能保证获得连续空间中的全局最优解。

    \item \textbf{候选数量采用稳定前缀截断。}

    超过数量上限时，
    当前实现只保留排序后的前$N_{\max}$个候选，
    未根据性能差距和参数覆盖性进行代表性采样。

    \item \textbf{不包含经验学习。}

    当前模板优先级和参数空间不会根据历史评估结果自动更新。

    \item \textbf{不在生成阶段判断SLA。}

    静态过滤只检查能力、参数、冲突和显式资源上限，
    不预测吞吐率、时延、语音质量和能耗。

    \item \textbf{初始生成与反馈生成相互独立。}

    初始候选由网格生成器产生，
    后续反馈候选由独立的局部反馈优化器生成，
    当前尚未形成统一的
    \texttt{SearchSpace}和
    \texttt{PolicyFeedback[]}生成接口。
\end{enumerate}

因此，网络策略生成完成的是：

\begin{equation}
\text{RAN目标与约束}
\longrightarrow
\text{动作模板}
\longrightarrow
\text{离散参数组合}
\longrightarrow
\text{静态可用候选策略}.
\label{eq:ran_requirement_to_candidate_policy}
\end{equation}

候选策略生成阶段只能说明策略在动作能力和静态规则层面可用，
不能说明策略实施后一定能够满足视频、语音和无线资源约束。

\section{候选策略评估}
\label{sec:candidate_policy_evaluation}

网络策略生成阶段根据业务意图转译结果、
RAN动作模板和离散参数空间，
生成了多个动作组合和参数配置不同的候选策略。
这些候选策略在业务保障效果、无线资源占用、
RAN能耗和潜在风险等方面可能存在明显差异，
因此不能直接将生成顺序靠前的策略作为最终推荐结果。

候选策略评估的任务是结合当前RAN状态，
计算每个候选策略实施后的网络性能，
检查其是否满足视频业务、语音业务和无线资源硬约束，
并在可行策略之间进行多目标评分和确定性排序。
对于未达到要求的策略，
系统还需要根据具体违约原因调整动作参数或动作组合，
形成新一轮局部候选策略。

候选策略评估过程可表示为：

\begin{equation}
\left(
X_{\mathrm{RAN}},
P_i
\right)
\xrightarrow{\mathrm{Performance\ Evaluation}}
Z_i
\xrightarrow{\mathrm{Constraint\ Check}}
H_i
\xrightarrow{\mathrm{Scoring}}
S_i
\xrightarrow{\mathrm{Feedback}}
\mathcal{N}(P_i),
\label{eq:candidate_policy_evaluation_process}
\end{equation}

其中：

\begin{itemize}
    \item $X_{\mathrm{RAN}}$表示当前单小区RAN状态；
    \item $P_i$表示第$i$个候选策略；
    \item $Z_i$表示策略实施后的性能评估结果；
    \item $H_i$表示硬约束检查结果；
    \item $S_i$表示多目标综合评分；
    \item $\mathcal{N}(P_i)$表示根据评估反馈生成的局部策略邻域。
\end{itemize}

与策略生成阶段只进行静态能力和冲突检查不同，
候选策略评估需要回答以下问题：

\begin{enumerate}[label=\arabic*)]
    \item 策略实施后，目标视频用户的吞吐率和时延是否满足要求；
    \item 语音业务质量是否保持在允许范围内；
    \item PRB利用率是否超过资源上限；
    \item 策略对RAN能耗、普通用户和系统风险产生什么影响；
    \item 哪些候选策略满足全部硬约束；
    \item 可行策略之间应如何评分和排序；
    \item 当策略不可行或仍有优化空间时，应如何生成下一轮候选策略。
\end{enumerate}


\subsection*{问题定义与总体方法}

从理论上看，
候选策略效果可以通过多种方式获得，例如：

\begin{itemize}
    \item 基于协议和无线传播模型的链路级仿真；
    \item 面向小区和多用户调度的系统级仿真；
    \item 基于历史数据训练的性能预测模型；
    \item 采用代理模型的快速策略评估；
    \item 在网络数字孪生环境中进行动态策略推演；
    \item 在真实网络沙箱或试验网中进行验证。
\end{itemize}

若采用通用性能评估器，
候选策略的性能结果可表示为：

\begin{equation}
Z_i
=
f_{\mathrm{eval}}
\left(
X_{\mathrm{RAN}},
P_i,
\boldsymbol{\theta}_{\mathrm{model}}
\right),
\label{eq:general_policy_evaluator}
\end{equation}

其中，
$\boldsymbol{\theta}_{\mathrm{model}}$
表示性能评估模型的参数。

理论上的高精度评估器需要考虑：

\begin{itemize}
    \item 无线传播和信道质量；
    \item 小区内多用户调度过程；
    \item PRB分配和缓存队列；
    \item HARQ、BLER和重传；
    \item 用户移动和小区间干扰；
    \item 业务流到达过程；
    \item 语音和视频业务模型；
    \item 基站射频、基带和载波能耗。
\end{itemize}

本文当前研究重点是验证从候选策略到评分、
排序和反馈优化的完整闭环，
尚未构建高精度网络数字孪生或协议级仿真环境。
因此，当前原型采用：

\begin{equation}
\boxed{
\text{简化RAN解析性能模型}
+
\text{硬约束检查}
+
\text{多目标加权评分}
+
\text{局部邻域反馈}
}
\label{eq:implemented_policy_evaluation_method}
\end{equation}

实现候选策略评估。

当前性能评估模型版本为：

\begin{equation}
ModelVersion
=
\texttt{simplified-ran-v1}.
\label{eq:performance_model_version}
\end{equation}

该模型通过参数化代数关系描述策略动作对吞吐率、
时延、BLER、语音质量、PRB利用率、
RAN能耗、普通用户性能和风险的相对影响。

需要说明的是，
该模型不包含真实协议栈、无线传播、队列调度和干扰过程，
只能用于验证相对趋势和决策链路，
不能解释为真实RAN性能预测。


\subsection*{简化RAN性能评估模型}

设候选策略的四个主要动作参数为：

\begin{equation}
\boldsymbol{\theta}_i
=
\left[
r_i,
w_i,
v_i,
e_i
\right],
\label{eq:policy_action_parameter_vector}
\end{equation}

其中：

\begin{itemize}
    \item $r_i$表示视频业务上行PRB预留比例；
    \item $w_i$表示视频业务调度权重；
    \item $v_i$表示语音业务最低PRB保护比例；
    \item $e_i$表示RAN节能控制强度。
\end{itemize}

若某个候选策略不包含对应动作，
性能模型使用配置中的中性参数：

\begin{equation}
r=0,
\qquad
w=1,
\qquad
v=0,
\qquad
e=0.
\label{eq:neutral_action_parameters}
\end{equation}

这样可以在统一公式中同时评估基线策略、
单动作策略和多动作联合策略。

当前RAN状态可以抽象表示为：

\begin{equation}
X_{\mathrm{RAN}}
=
\left\{
u_0,
E_0,
Q_0,
\left(
T_u^0,
D_u^0,
B_u^0,
\gamma_u
\right)_{u\in U_{\mathrm{video}}}
\right\},
\label{eq:ran_state_for_evaluation}
\end{equation}

其中：

\begin{itemize}
    \item $u_0$表示基线PRB利用率；
    \item $E_0$表示基线RAN能耗；
    \item $Q_0$表示基线语音质量；
    \item $T_u^0$表示视频用户$u$的基线上行吞吐率；
    \item $D_u^0$表示视频用户$u$的基线上行时延；
    \item $B_u^0$表示视频用户$u$的基线BLER；
    \item $\gamma_u$表示视频用户$u$的SINR或信道质量输入。
\end{itemize}


\subsubsection*{有效容量与PRB利用率}

节能动作可能降低可用系统容量。
当前模型定义有效容量因子：

\begin{equation}
\phi
=
\operatorname{clip}
\left(
1-k_E e,
\phi_{\min},
1
\right),
\label{eq:effective_capacity_factor}
\end{equation}

其中，
$k_E$表示节能强度对有效容量的影响系数，
$\phi_{\min}$表示容量下界。

策略实施后的PRB利用率为：

\begin{equation}
u'
=
\operatorname{clip}
\left(
\frac{
u_0+k_r r+k_v v
}{
\phi
},
0,
1
\right),
\label{eq:post_policy_prb_utilization}
\end{equation}

其中：

\begin{itemize}
    \item $k_r$描述视频PRB预留对资源利用率的影响；
    \item $k_v$描述语音资源保护对资源利用率的影响。
\end{itemize}

当策略后利用率超过拥塞门限
$u_{\mathrm{cong}}$时，
定义拥塞量：

\begin{equation}
o
=
\max
\left(
0,
u'-u_{\mathrm{cong}}
\right).
\label{eq:congestion_amount}
\end{equation}

拥塞量将进一步影响视频吞吐率、
时延、BLER、语音质量和风险。


\subsubsection*{视频用户吞吐率}

模型首先根据用户SINR计算信道质量影响因子：

\begin{equation}
f_u^{\mathrm{SINR}}
=
\operatorname{clip}
\left(
1+
k_s
\left(
\gamma_u-\gamma_{\mathrm{ref}}
\right),
f_{\min},
f_{\max}
\right),
\label{eq:sinr_factor}
\end{equation}

其中，
$\gamma_{\mathrm{ref}}$表示参考SINR，
$k_s$表示SINR变化对吞吐率的线性影响系数。

视频用户$u$策略后的上行吞吐率为：

\begin{equation}
T_u'
=
\operatorname{clip}
\left[
T_u^0
\left(
1+\alpha_r r+\alpha_w(w-1)
\right)
\phi
\frac{
f_u^{\mathrm{SINR}}
}{
1+\alpha_c o
},
T_{\min},
T_{\max}
\right],
\label{eq:post_policy_video_throughput}
\end{equation}

其中：

\begin{itemize}
    \item $\alpha_r$表示PRB预留的吞吐率增益系数；
    \item $\alpha_w$表示调度权重的吞吐率增益系数；
    \item $\alpha_c$表示拥塞对吞吐率的抑制系数；
    \item $T_{\min}$和$T_{\max}$表示吞吐率输出边界。
\end{itemize}

从式\eqref{eq:post_policy_video_throughput}可以看出，
PRB预留和调度权重能够提高视频吞吐率，
节能容量折减和系统拥塞则会降低吞吐率。


\subsubsection*{视频用户时延}

视频用户$u$策略后的上行时延为：

\begin{equation}
D_u'
=
\operatorname{clip}
\left[
D_u^0
\frac{
1+
\beta_u
\max(0,u'-u_0)
+
\beta_e e
}{
1+
\beta_r r+
\beta_w(w-1)
},
D_{\min},
D_{\max}
\right],
\label{eq:post_policy_video_delay}
\end{equation}

其中：

\begin{itemize}
    \item $\beta_u$表示资源利用率增长对时延的影响；
    \item $\beta_e$表示节能动作对时延的影响；
    \item $\beta_r$表示PRB预留的时延改善系数；
    \item $\beta_w$表示调度权重的时延改善系数；
    \item $D_{\min}$和$D_{\max}$表示时延输出边界。
\end{itemize}

式\eqref{eq:post_policy_video_delay}表明，
PRB预留和调度权重能够降低视频时延，
而资源利用率增加和节能动作可能提高时延。


\subsubsection*{视频用户BLER}

策略后的用户级BLER为：

\begin{equation}
B_u'
=
\operatorname{clip}
\left[
B_u^0
+
\chi_u o
-
\chi_r r
-
\chi_w(w-1),
0,
1
\right],
\label{eq:post_policy_video_bler}
\end{equation}

其中：

\begin{itemize}
    \item $\chi_u$表示拥塞对BLER的恶化作用；
    \item $\chi_r$表示PRB预留对BLER的改善作用；
    \item $\chi_w$表示调度权重对BLER的改善作用。
\end{itemize}

当前模型将BLER作为软风险和业务评分指标，
而不是独立硬约束。


\subsubsection*{用户级指标聚合}

对所有目标视频用户，
系统计算以下聚合指标：

\begin{itemize}
    \item 平均上行吞吐率；
    \item 低分位上行吞吐率；
    \item 平均上行时延；
    \item 高分位上行时延；
    \item 平均BLER；
    \item 满足单用户吞吐率门限的用户比例。
\end{itemize}

平均吞吐率为：

\begin{equation}
\overline{T}
=
\frac{1}
{|U_{\mathrm{video}}|}
\sum_{u\in U_{\mathrm{video}}}
T_u'.
\label{eq:average_video_throughput}
\end{equation}

低分位吞吐率为：

\begin{equation}
T_{\mathrm{low}}
=
Q_{q_T}
\left(
\left\{
T_u'
\mid
u\in U_{\mathrm{video}}
\right\}
\right),
\label{eq:low_quantile_video_throughput}
\end{equation}

其中，
$q_T$由配置文件给出，
当前正式实验中对应P05统计值。

高分位时延为：

\begin{equation}
D_{\mathrm{high}}
=
Q_{q_D}
\left(
\left\{
D_u'
\mid
u\in U_{\mathrm{video}}
\right\}
\right),
\label{eq:high_quantile_video_delay}
\end{equation}

其中，
$q_D$由配置文件给出，
当前正式实验中对应P95统计值。

平均BLER为：

\begin{equation}
\overline{B}
=
\frac{1}
{|U_{\mathrm{video}}|}
\sum_{u\in U_{\mathrm{video}}}
B_u'.
\label{eq:average_video_bler}
\end{equation}

视频用户达标比例为：

\begin{equation}
\rho_T
=
\frac{
\sum_{u\in U_{\mathrm{video}}}
\mathbb{I}
\left[
T_u'\geq T_{\mathrm{th}}
\right]
}{
|U_{\mathrm{video}}|
}.
\label{eq:video_user_satisfaction_ratio_eval}
\end{equation}


\subsubsection*{语音质量}

策略实施后的语音质量定义为：

\begin{equation}
Q'
=
\operatorname{clip}
\left(
Q_0
+
\delta_v v
-
\delta_r r
-
\delta_e e
-
\delta_u o,
0,
1
\right),
\label{eq:post_policy_voice_quality}
\end{equation}

其中：

\begin{itemize}
    \item $\delta_v$表示语音资源保护带来的质量增益；
    \item $\delta_r$表示视频PRB预留对语音资源竞争的影响；
    \item $\delta_e$表示节能动作对语音质量的影响；
    \item $\delta_u$表示拥塞对语音质量的影响。
\end{itemize}

语音质量相对基线的下降为：

\begin{equation}
\Delta Q_{\mathrm{voice}}
=
\operatorname{clip}
\left(
\max
\left(
0,
Q_0-Q'
\right),
0,
1
\right).
\label{eq:voice_quality_degradation_eval}
\end{equation}


\subsubsection*{RAN能耗}

策略后的RAN能耗为：

\begin{equation}
E'
=
\operatorname{clip}
\left[
E_0
\left(
1-\eta_e e
+
\eta_u
\max(0,u'-u_0)
+
\eta_a
N_{\mathrm{highcost}}
\right),
E_{\min},
E_{\max}
\right],
\label{eq:post_policy_ran_energy}
\end{equation}

其中：

\begin{itemize}
    \item $\eta_e$表示节能动作带来的能耗降低；
    \item $\eta_u$表示资源利用率增长带来的额外能耗；
    \item $\eta_a$表示高开销动作带来的附加能耗；
    \item $N_{\mathrm{highcost}}$表示策略中高开销动作数量。
\end{itemize}

高开销动作不是在性能模型中固定写死，
而是由动作库中的
\texttt{isHighCostAction}
属性确定。


\subsubsection*{普通用户性能影响}

为了描述业务保障动作对非目标普通用户的影响，
模型定义普通用户性能下降：

\begin{equation}
\Delta Q_{\mathrm{normal}}
=
\operatorname{clip}
\left[
\zeta_r r
+
\zeta_w(w-1)
+
\zeta_v v
+
\zeta_e e
+
\zeta_u
\max(0,u'-u_0),
0,
1
\right].
\label{eq:normal_user_degradation}
\end{equation}

该指标不是硬约束，
而是作为公平性风险进入综合风险评分。


\subsubsection*{风险建模}

当前模型从以下方面构造风险：

\begin{itemize}
    \item PRB过载风险；
    \item 语音质量风险；
    \item 激进节能风险；
    \item 动作数量和控制复杂度风险；
    \item 普通用户性能下降风险；
    \item 视频BLER恶化风险。
\end{itemize}

PRB过载风险为：

\begin{equation}
R_u
=
\operatorname{clip}
\left(
\frac{
u'-u_{\mathrm{cong}}
}{
1-u_{\mathrm{cong}}
},
0,
1
\right).
\label{eq:overload_risk}
\end{equation}

语音绝对质量风险可表示为：

\begin{equation}
R_{q,a}
=
\operatorname{clip}
\left(
\frac{
Q_{\min}+m-Q'
}{
m
},
0,
1
\right),
\label{eq:voice_absolute_risk}
\end{equation}

其中，
$m$表示语音质量风险归一化裕量。

语音相对下降风险为：

\begin{equation}
R_{q,b}
=
\operatorname{clip}
\left(
\frac{
\Delta Q_{\mathrm{voice}}
}{
\max
\left(
\Delta Q_{\max},
\epsilon
\right)
},
0,
1
\right).
\label{eq:voice_relative_risk}
\end{equation}

最终语音风险取两者较大值：

\begin{equation}
R_q
=
\max
\left(
R_{q,a},
R_{q,b}
\right).
\label{eq:voice_risk}
\end{equation}

节能风险为：

\begin{equation}
R_e
=
\operatorname{clip}
\left(
\frac{
e
}{
\max(e_{\mathrm{configured}},\epsilon)
},
0,
1
\right).
\label{eq:energy_risk}
\end{equation}

动作复杂度风险为：

\begin{equation}
R_a
=
\operatorname{clip}
\left(
\frac{
N_{\mathrm{active}}
}{
N_{\max}
},
0,
1
\right).
\label{eq:action_count_risk}
\end{equation}

普通用户公平性风险为：

\begin{equation}
R_f
=
\operatorname{clip}
\left(
\frac{
\Delta Q_{\mathrm{normal}}
-
\theta_f
}{
1-\theta_f
},
0,
1
\right),
\label{eq:fairness_risk}
\end{equation}

其中，
$\theta_f$表示普通用户性能下降软门限。

BLER风险为：

\begin{equation}
R_b
=
\operatorname{clip}
\left(
\frac{
\max
\left(
0,
\overline{B}'-\overline{B}_0
\right)
}{
s_b
},
0,
1
\right),
\label{eq:bler_risk}
\end{equation}

其中，
$s_b$表示BLER风险尺度。

综合风险为：

\begin{equation}
R
=
\operatorname{clip}
\left(
\omega_uR_u+
\omega_qR_q+
\omega_eR_e+
\omega_aR_a+
\omega_fR_f+
\omega_bR_b,
0,
1
\right),
\label{eq:total_policy_risk}
\end{equation}

其中，
各风险权重均由配置文件给出。


\subsection*{硬约束检查与多目标评分}

性能模型完成策略效果计算后，
系统首先执行硬约束检查，
然后只对满足全部硬约束的策略进行有效排序。

硬约束检查器遍历策略中携带的全部约束，
不会在发现第一项违约后提前停止，
而是保留所有违约原因，
为后续反馈优化提供诊断信息。

设第$j$项硬约束满足函数为：

\begin{equation}
c_j(P_i)
=
\begin{cases}
1,
& \text{策略$P_i$满足第$j$项约束},\\
0,
& \text{策略$P_i$违反第$j$项约束}.
\end{cases}
\label{eq:single_hard_constraint}
\end{equation}

候选策略的总体可行性为：

\begin{equation}
H(P_i)
=
\prod_{j=1}^{J}
c_j(P_i).
\label{eq:hard_constraint_feasibility}
\end{equation}

当全部硬约束均满足时：

\begin{equation}
H(P_i)=1;
\end{equation}

否则：

\begin{equation}
H(P_i)=0.
\end{equation}

当前演唱会场景实际检查以下硬约束：

\begin{equation}
\left\{
\begin{aligned}
&\rho_T\geq0.95,\\
&D_{\mathrm{video}}^{P95}\leq30~\mathrm{ms},\\
&Q'\geq0.95,\\
&\Delta Q_{\mathrm{voice}}\leq0.02,\\
&u'\leq0.90.
\end{aligned}
\right.
\label{eq:hard_constraint_set_evaluation}
\end{equation}

对于不可行策略，
本文采用硬约束不可补偿原则：

\begin{equation}
H(P_i)=0
\quad\Longrightarrow\quad
Score(P_i)=0,
\qquad
Rank(P_i)=\varnothing.
\label{eq:infeasible_policy_score_rule}
\end{equation}

这意味着节能收益、低风险或低资源开销，
不能补偿视频SLA、语音保护或PRB上限违约。

不可行策略仍保留性能结果、
归一化指标和各子评分，
用于分析违约原因和生成反馈，
但其总分固定为0且不进入可行策略排名。


\subsubsection*{指标归一化}

不同评价指标的量纲和优化方向不同。
本文采用固定上下界归一化，
而不是使用当前候选集合的最小值和最大值。

对于增益型指标，
定义：

\begin{equation}
N_g(x)
=
\operatorname{clip}
\left(
\frac{
x-x_{\min}
}{
x_{\max}-x_{\min}
},
0,
1
\right).
\label{eq:gain_normalization}
\end{equation}

对于成本型指标，
定义：

\begin{equation}
N_c(x)
=
1-N_g(x).
\label{eq:cost_normalization}
\end{equation}

归一化上下界均由配置文件提供。
采用固定边界能够保证不同实验批次之间的评分口径保持一致。


\subsubsection*{业务保障评分}

当前业务保障评分综合考虑：

\begin{itemize}
    \item 视频用户达标比例；
    \item 视频P05吞吐率；
    \item 视频P95时延；
    \item 视频平均BLER；
    \item 语音质量。
\end{itemize}

其计算公式为：

\begin{equation}
\begin{aligned}
S_{\mathrm{service}}
={}&
0.35N_g(\rho_T)
+
0.20N_g(T_{P05})
+
0.20N_c(D_{P95})
\\
&+
0.10N_c(\overline{B})
+
0.15N_g(Q').
\end{aligned}
\label{eq:service_score}
\end{equation}


\subsubsection*{节能评分}

节能收益定义为策略相对于基线状态的能耗下降比例：

\begin{equation}
G_E
=
\frac{
E_0-E'
}{
E_0
}.
\label{eq:energy_saving_gain}
\end{equation}

节能评分为：

\begin{equation}
S_{\mathrm{energy}}
=
N_g(G_E).
\label{eq:energy_score}
\end{equation}


\subsubsection*{资源评分}

资源评分同时考虑策略后的PRB利用率和动作数量：

\begin{equation}
S_{\mathrm{resource}}
=
0.80N_c(u')
+
0.20N_c(N_{\mathrm{actions}}).
\label{eq:resource_score}
\end{equation}

该评分倾向于选择资源占用较低、
动作组合较简单的候选策略。


\subsubsection*{风险评分}

风险评分为综合风险的成本型归一化结果：

\begin{equation}
S_{\mathrm{risk}}
=
N_c(R).
\label{eq:risk_score}
\end{equation}

风险越低，
$S_{\mathrm{risk}}$越高。


\subsubsection*{综合评分}

对满足全部硬约束的候选策略，
综合评分为：

\begin{equation}
Score(P_i)
=
0.60S_{\mathrm{service}}
+
0.15S_{\mathrm{energy}}
+
0.15S_{\mathrm{resource}}
+
0.10S_{\mathrm{risk}}.
\label{eq:total_policy_score}
\end{equation}

更一般地，
可以表示为：

\begin{equation}
Score(P_i)
=
\sum_{k=1}^{4}
\lambda_k
S_k(P_i),
\label{eq:general_weighted_policy_score}
\end{equation}

满足：

\begin{equation}
\sum_{k=1}^{4}
\lambda_k
=
1,
\qquad
\lambda_k\geq0.
\label{eq:score_weight_constraint}
\end{equation}

当前默认权重为：

\begin{equation}
\boldsymbol{\lambda}
=
\left[
0.60,
0.15,
0.15,
0.10
\right],
\label{eq:default_score_weights}
\end{equation}

分别对应业务、节能、资源和风险评分。

所有权重均从配置文件读取，
并由数据模型校验权重之和为1。


\subsubsection*{候选策略排序}

当前源码对所有可行策略首先按综合评分降序排序。
需要注意，
当前实现不是“先比较业务评分，再比较综合评分”，
而是直接以综合评分作为第一排序依据。

策略排序主键为：

\begin{equation}
K(P_i)
=
\left(
-Score(P_i),
u_i',
R_i,
E_i',
N_{\mathrm{actions},i},
ID_i
\right).
\label{eq:policy_ranking_key}
\end{equation}

当多个候选策略综合评分相同或在数值舍入后相同，
依次选择：

\begin{enumerate}[label=\arabic*)]
    \item PRB利用率更低的策略；
    \item 综合风险更低的策略；
    \item RAN能耗更低的策略；
    \item 动作数量更少的策略；
    \item 策略标识字典序更小的策略。
\end{enumerate}

这种排序方式能够保证结果完全确定，
即使多个策略评分相同，
也能获得稳定、可复现的排名。

不可行策略统一放在可行策略之后，
且不赋予排名。


\subsection*{反馈优化与终止机制}

候选策略评估不仅用于选出当前最优策略，
还需要将性能结果和违约原因反馈给策略生成过程。

从理论上看，
反馈优化可以采用：

\begin{itemize}
    \item 梯度优化；
    \item 数学规划；
    \item 模型预测控制；
    \item 贝叶斯优化；
    \item 强化学习；
    \item 多目标进化算法；
    \item 基于数字孪生的滚动优化。
\end{itemize}

本文当前采用：

\begin{equation}
\boxed{
\text{单父策略}
+
\text{相邻离散档位}
+
\text{单操作局部邻域搜索}
}
\label{eq:implemented_feedback_method}
\end{equation}

完成候选策略改进。

设第$n$轮选中的父策略为
$P_{\mathrm{parent}}^{(n)}$，
反馈优化器根据该策略的约束违约和性能薄弱项，
生成下一轮候选集合：

\begin{equation}
\mathcal{P}^{(n+1)}
=
\mathcal{N}
\left(
P_{\mathrm{parent}}^{(n)},
F^{(n)}
\right),
\label{eq:next_round_policy_neighborhood}
\end{equation}

其中，
$F^{(n)}$表示第$n$轮评估反馈。


\subsubsection*{父策略选择}

当当前轮存在可行策略时，
系统选择排名第1的策略作为反馈父策略：

\begin{equation}
P_{\mathrm{parent}}^{(n)}
=
\arg\min_{P_i}
Rank(P_i).
\label{eq:feasible_parent_selection}
\end{equation}

当当前轮所有策略均不可行时，
系统根据诊断性子评分选择相对较优的父策略，
依次偏好：

\begin{itemize}
    \item 较高的业务保障评分；
    \item 较高的资源评分；
    \item 较高的风险评分；
    \item 较低的PRB利用率；
    \item 较小的策略标识。
\end{itemize}

这种方式使反馈优化即使在无可行策略时，
仍能从相对接近业务目标的候选开始改进。


\subsubsection*{相邻参数调整}

当前动作参数均来自有序离散集合。
参数“增加一档”定义为选择当前值之后的第一个更大允许值，
参数“降低一档”定义为选择当前值之前的第一个更小允许值。

可以表示为：

\begin{equation}
\theta^{(n+1)}
=
\operatorname{NextAllowed}
\left(
\theta^{(n)},
d
\right),
\label{eq:adjacent_parameter_update}
\end{equation}

其中：

\begin{equation}
d
\in
\left\{
+1,-1
\right\}
\end{equation}

分别表示增加一档或降低一档。

若父策略中不存在某类动作，
则优化器从性能模型中对应的中性值出发，
尝试增加该动作。

当前主要反馈规则如表
\ref{tab:implemented_feedback_rules}
所示。

\begin{table}[H]
    \centering
    \caption{当前实现的候选策略反馈规则}
    \label{tab:implemented_feedback_rules}
    \small
    \renewcommand{\arraystretch}{1.25}
    \begin{tabularx}{\textwidth}{
        >{\raggedright\arraybackslash}p{0.34\textwidth}
        >{\raggedright\arraybackslash}X
    }
        \toprule
        评估结果或触发条件 & 局部调整操作 \\
        \midrule

        视频用户达标比例不足或P05吞吐率过低
        &
        增加PRB预留比例；
        增加视频调度权重
        \\

        视频P95时延超过门限
        &
        增加视频调度权重；
        增加PRB预留比例
        \\

        语音绝对质量或相对下降违约
        &
        增加语音最低PRB比例；
        降低视频PRB预留；
        降低节能强度
        \\

        策略后PRB利用率过高
        &
        降低视频PRB预留；
        删除动作库中标记的高开销动作
        \\

        策略可行但RAN能耗高于基线
        &
        增加节能强度；
        降低视频PRB预留
        \\

        普通用户性能下降超过软门限
        &
        降低视频PRB预留；
        降低视频调度权重
        \\

        视频BLER风险增加
        &
        降低节能强度；
        降低视频PRB预留
        \\

        无可行策略且包含节能软目标
        &
        在配置允许时删除节能动作，
        并移除RAN节能优化目标
        \\

        \bottomrule
    \end{tabularx}
\end{table}

需要说明的是，
同一个评估反馈可能产生多个调整建议，
但当前优化器不会把多个调整一次性合并到同一新策略中。

例如，
当视频吞吐率不足时，
系统分别生成：

\begin{equation}
P_{r+}
=
P_{\mathrm{parent}}
\oplus
\left(
r\uparrow
\right),
\end{equation}

和：

\begin{equation}
P_{w+}
=
P_{\mathrm{parent}}
\oplus
\left(
w\uparrow
\right).
\end{equation}

每个局部操作独立形成一个新候选，
从而便于分析单个参数变化的影响。


\subsubsection*{动作增加与删除}

反馈优化不仅可以调整已有动作参数，
也可以增加或删除动作。

当策略缺少某类必要动作时，
优化器可以从中性值出发，
增加一个对应动作。

例如：

\begin{equation}
P'
=
P
\cup
\left\{
a_{\mathrm{voice}}
\right\}
\end{equation}

表示增加语音资源保护动作。

当PRB过载或节能软目标导致不可行时，
系统可以删除某些动作：

\begin{equation}
P'
=
P
\setminus
\left\{
a_{\mathrm{energy}}
\right\}.
\end{equation}

但语音保护动作在发生语音违约时不会被删除。

所有反馈候选在生成后，
仍需重新经过静态约束检查，
不符合能力、参数或冲突规则的候选直接丢弃。


\subsubsection*{反馈候选标识与去重}

反馈候选保存：

\begin{itemize}
    \item 父策略标识；
    \item 下一轮生成轮次；
    \item 反馈原因；
    \item 局部调整操作。
\end{itemize}

反馈候选的策略标识基于规范化策略内容的哈希生成：

\begin{equation}
PolicyId
=
\texttt{RAN-POLICY-R}
\Vert
n
\Vert
\texttt{-FB-}
\Vert
Hash(P)_{\mathrm{prefix}}.
\label{eq:feedback_policy_id}
\end{equation}

该机制使相同内容的反馈候选获得稳定标识。

优化器会排除自身已经生成过的策略签名，
避免在连续反馈中反复生成同一局部变体。

需要指出的是，
当前端到端流水线没有将每一轮所有已评估策略的完整签名集合
全部传给反馈优化器，
因此还不能严格保证不会重新评估某个非父初始策略。
这是当前实现需要进一步完善的地方。


\subsubsection*{反馈开关}

当前原型提供：

\begin{equation}
feedbackEnabled
\in
\left\{
\texttt{true},
\texttt{false}
\right\}
\end{equation}

用于控制是否执行反馈优化。

当：

\begin{equation}
feedbackEnabled
=
\texttt{false}
\end{equation}

时，
系统仍完整执行：

\begin{itemize}
    \item 第0轮候选策略生成；
    \item 性能评估；
    \item 硬约束检查；
    \item 多目标评分；
    \item 策略排序。
\end{itemize}

随后立即停止，
不调用反馈优化器，
不生成后续轮次。

此时输出记录：

\begin{equation}
feedbackIterations=0,
\end{equation}

\begin{equation}
generationRounds=1,
\end{equation}

\begin{equation}
terminationReason
=
\texttt{feedback\_disabled}.
\end{equation}

其中，
\texttt{generationRounds=1}
表示实际执行了一个生成轮，
该轮编号为0。


\subsubsection*{终止条件}

当反馈开启时，
系统在每轮评估后检查是否满足终止条件。

主要终止条件包括：

\begin{enumerate}[label=\arabic*)]
    \item 已存在满足全部硬约束，
    且总分达到配置门限的策略；
    \item 最近若干轮最优分提升均低于最小改进阈值；
    \item 达到最大反馈迭代次数；
    \item 当前父策略无法生成新的有效邻域候选。
\end{enumerate}

评分门限停止条件为：

\begin{equation}
\max_{P_i\in\mathcal{P}^{(n)}}
Score(P_i)
\geq
Score_{\mathrm{th}}.
\label{eq:score_threshold_termination}
\end{equation}

收敛条件可以表示为：

\begin{equation}
\left|
Score_{\max}^{(n)}
-
Score_{\max}^{(n-1)}
\right|
<
\varepsilon_{\mathrm{score}},
\label{eq:score_convergence_condition}
\end{equation}

并要求该条件连续满足配置轮数。

最大反馈次数条件为：

\begin{equation}
n_{\mathrm{feedback}}
\geq
N_{\mathrm{feedback}}^{\max}.
\label{eq:maximum_feedback_iterations}
\end{equation}

整个反馈过程始终保持：

\begin{equation}
\texttt{allowRelaxHardConstraints}
=
\texttt{false}.
\label{eq:no_hard_constraint_relaxation}
\end{equation}

即使不存在可行策略，
系统也不会降低视频、语音或PRB硬约束。
当前唯一允许的目标调整是，
在配置允许时取消节能软目标。


\subsection*{评估输出与能力边界}

候选策略评估最终输出：

\begin{equation}
O_{\mathrm{evaluation}}
=
\left\{
P_i,
Z_i,
H_i,
S_i,
Rank_i,
F_i
\right\}_{i=1}^{M},
\label{eq:evaluation_output_model}
\end{equation}

其中：

\begin{itemize}
    \item $P_i$表示候选策略；
    \item $Z_i$表示性能评估结果；
    \item $H_i$表示硬约束检查结果；
    \item $S_i$表示子评分和综合评分；
    \item $Rank_i$表示可行策略排名；
    \item $F_i$表示反馈诊断和改进建议。
\end{itemize}

一个简化的评估输出示例如下：

\begin{lstlisting}[language={},
caption={候选策略评估结果示例},
label={lst:candidate_policy_evaluation_example}]
{
  "intentId": "INTENT-CONCERT-UL-001",
  "generationRound": 0,

  "candidatePolicies": [
    {
      "policyId": "RAN-POLICY-R0-0016",

      "performance": {
        "videoUserSatisfactionRatio": 0.97,
        "videoThroughputP05Mbps": 22.28,
        "videoDelayP95Ms": 29.18,
        "videoUplinkBlerMean": 0.08,
        "voiceQualityScore": 0.95,
        "voiceQualityDegradation": 0.01,
        "prbUtilization": 0.84,
        "ranEnergy": 0.91,
        "normalUserDegradation": 0.06,
        "riskScore": 0.24
      },

      "constraintCheck": {
        "feasible": true,
        "violations": []
      },

      "scores": {
        "serviceScore": 0.71,
        "energyScore": 0.45,
        "resourceScore": 0.39,
        "riskScore": 0.76,
        "totalScore": 0.61
      },

      "rank": 1
    },

    {
      "policyId": "RAN-POLICY-R0-0021",

      "performance": {
        "videoUserSatisfactionRatio": 0.91,
        "videoThroughputP05Mbps": 18.70,
        "videoDelayP95Ms": 34.20,
        "voiceQualityScore": 0.96,
        "voiceQualityDegradation": 0.01,
        "prbUtilization": 0.79
      },

      "constraintCheck": {
        "feasible": false,
        "violations": [
          "video_user_satisfaction_ratio",
          "video_delay_p95"
        ]
      },

      "scores": {
        "serviceScore": 0.54,
        "energyScore": 0.66,
        "resourceScore": 0.52,
        "riskScore": 0.72,
        "totalScore": 0.0
      },

      "rank": null,

      "feedback": [
        "increase_reserved_prb_ratio",
        "increase_video_scheduling_weight"
      ]
    }
  ],

  "bestPolicyId": "RAN-POLICY-R0-0016",
  "nextAction": "feedback_or_stop"
}
\end{lstlisting}

需要说明的是，
上述数据仅用于展示输出结构，
正式结果应以实际仿真实验产物为准。

当前候选策略评估方法具有以下优点：

\begin{itemize}
    \item 采用统一性能评估接口，
    便于未来替换为网络数字孪生评估器；
    \item 对视频用户进行逐用户计算，
    并输出P05、P95和达标比例等统计指标；
    \item 硬约束检查不短路，
    能够保留全部违约原因；
    \item 不可行策略不允许由节能或资源收益补偿；
    \item 评分边界和权重均来自版本化配置；
    \item 排序和并列规则完全确定；
    \item 反馈候选保留父子关系和生成原因；
    \item 整个过程可重复执行并进行自动化测试。
\end{itemize}

但当前实现仍存在以下能力边界：

\begin{enumerate}[label=\arabic*)]
    \item \textbf{性能模型为简化解析模型。}

    当前模型不包含真实无线传播、
    协议栈、队列调度、HARQ和小区间干扰，
    只能描述预设动作空间中的相对趋势。

    \item \textbf{模型系数属于仿真假设。}

    吞吐率、时延、语音和能耗公式中的系数，
    均来自版本化配置，
    不是由现网数据标定得到的真实参数。

    \item \textbf{不包含真正的随机模式。}

    当前模型记录随机种子，
    但性能计算不执行随机采样，
    不同种子下结果保持一致。

    \item \textbf{评分采用固定加权和。}

    当前以综合总分作为第一排序键，
    不使用词典序多目标优化、
    $\varepsilon$约束法或在线权重调整。

    \item \textbf{反馈只搜索局部邻域。}

    当前每轮只选择一个父策略，
    并对单个参数或动作执行一次局部调整，
    无法保证获得全局最优策略。

    \item \textbf{没有联合多操作搜索。}

    同一反馈中的多个操作不会组合为一个新候选，
    因而可能错过需要多参数同时变化的改进方向。

    \item \textbf{不更新长期策略知识。}

    当前评估结果不会自动更新模板优先级、
    历史成功率或动作经验库。

    \item \textbf{动作空间不自动扩展。}

    配置中虽然保留动作空间扩展开关，
    当前算法并未实际读取和使用该功能。

    \item \textbf{尚未接入高精度网络数字孪生。}

    当前只通过统一接口预留未来替换能力，
    第四章将进一步研究面向策略验证的网络数字孪生建模与编排。
\end{enumerate}

候选策略评估最终完成：

\begin{equation}
\text{候选策略}
\longrightarrow
\text{性能结果}
\longrightarrow
\text{硬约束判定}
\longrightarrow
\text{多目标排序}
\longrightarrow
\text{局部反馈改进}.
\label{eq:candidate_to_ranked_policy}
\end{equation}

至此，
本文构建了从自然语言业务需求、
结构化意图、RAN需求、
候选策略生成到策略评估与反馈优化的完整闭环。

\section{仿真实现与结果分析}
\label{sec:simulation_results}

前文分别介绍了业务意图抽取、业务意图转译、
候选RAN策略生成以及候选策略评估与反馈优化方法。
为验证上述方法能否形成完整、稳定和可复现的决策闭环，
本文基于当前单小区RAN仿真原型开展正式批量实验。

实验主要验证以下内容：

\begin{enumerate}[label=\arabic*)]
    \item 自然语言业务需求能否依次完成意图抽取、
    意图转译、策略生成、性能评估、约束检查和策略推荐；

    \item 在不同网络负载条件下，
    候选策略的可行空间和性能表现如何变化；

    \item 开启反馈优化后，
    是否能够发现新的可行策略或提升最优策略性能；

    \item 改变业务、节能、资源和风险评分权重后，
    策略排序是否保持稳定；

    \item 简化性能模型中的输出裁剪是否会改变硬约束判定
    和最终推荐结果。
\end{enumerate}


\subsection*{仿真环境与实验设置}

当前仿真原型采用Python实现，
主要依赖Pydantic完成数据模型校验，
利用NumPy和Pandas进行性能计算与实验结果处理，
并使用Pytest完成单元测试和端到端测试。

当前实验采用的主要软件和模型版本为：

\begin{itemize}
    \item Python版本：3.11；
    \item 仿真原型版本：v0.1.1；
    \item 性能模型版本：\texttt{simplified-ran-v1}；
    \item 网络范围：单小区RAN；
    \item 策略动作空间：离散动作参数集合；
    \item 性能计算模式：确定性模式；
    \item 随机采样：未启用；
    \item 反馈优化：根据实验配置开启或关闭。
\end{itemize}

正式实验共包含31个实验实例，
分为基线实验、负载扫描实验、
反馈机制对比实验和评分权重敏感性实验。

实验规模如表
\ref{tab:formal_experiment_scale}
所示。

\begin{table}[H]
    \centering
    \caption{正式仿真实验规模}
    \label{tab:formal_experiment_scale}
    \small
    \renewcommand{\arraystretch}{1.25}
    \begin{tabular}{
        >{\raggedright\arraybackslash}p{0.30\textwidth}
        >{\centering\arraybackslash}p{0.14\textwidth}
        >{\centering\arraybackslash}p{0.18\textwidth}
        >{\centering\arraybackslash}p{0.18\textwidth}
    }
        \toprule
        实验类型 & 实例数 & 评估策略数 & 可行策略数 \\
        \midrule

        基线实验
        & 1
        & 109
        & 15
        \\

        负载扫描实验
        & 9
        & 900
        & 588
        \\

        反馈机制对比实验
        & 6
        & 627
        & 87
        \\

        评分权重敏感性实验
        & 15
        & 1500
        & 210
        \\

        \midrule
        合计
        & 31
        & 3136
        & 900
        \\

        \bottomrule
    \end{tabular}
\end{table}

全部31个实验实例均成功完成，
共评估3136个候选策略，
其中900个策略满足全部硬约束。

当前实验中的业务门限、
动作参数、评分权重和性能模型系数均为仿真假设，
其主要作用是验证方法流程和相对变化趋势，
不代表现网参数或标准规定。


\subsection*{端到端决策链路验证}

首先对完整意图驱动决策链路进行验证。

仿真输入采用如下自然语言业务需求：

\begin{quote}
“今晚演唱会期间，保障体育场内视频用户的上行体验，
尽量减少能耗，同时不影响语音业务。”
\end{quote}

系统依次完成：

\begin{equation}
\begin{aligned}
\text{Natural Language}
&\rightarrow
\text{Structured Intent}
\rightarrow
\text{Translated RAN Intent}
\\
&\rightarrow
\text{Candidate Policies}
\rightarrow
\text{Performance Results}
\\
&\rightarrow
\text{Constraint Check}
\rightarrow
\text{Policy Ranking}
\rightarrow
\text{Feedback}.
\end{aligned}
\label{eq:end_to_end_simulation_chain}
\end{equation}

在业务意图抽取阶段，
系统识别出：

\begin{itemize}
    \item 视频上行保障意图；
    \item RAN节能优化意图；
    \item 语音业务保护意图；
    \item 时间、地点、目标对象和约束关系；
    \item 准确时间、目标小区和SLA门限等待解析项。
\end{itemize}

在业务意图转译阶段，
静态事件库和场馆—小区映射库补全活动时间和目标小区，
SLA模板补全视频、语音和PRB约束，
转译状态由\texttt{blocked}转换为\texttt{ready}。

在候选策略生成阶段，
系统根据目标选择策略模板，
对PRB预留比例、视频调度权重、
语音资源保护比例和节能强度执行离散组合，
并通过静态资源约束和动作冲突规则过滤候选策略。

在候选策略评估阶段，
系统计算每个策略对应的：

\begin{itemize}
    \item 视频用户达标比例；
    \item 视频P05上行吞吐率；
    \item 视频P95上行时延；
    \item 平均BLER；
    \item 语音质量及其相对下降；
    \item PRB利用率；
    \item RAN能耗；
    \item 普通用户性能下降；
    \item 综合风险。
\end{itemize}

随后执行硬约束检查和多目标评分。
不可行策略总分置为0且不赋予排名，
可行策略按照综合评分和并列规则进行确定性排序。

当反馈开启时，
系统根据当前最优策略或最接近可行的策略，
生成相邻参数或动作变体，
并继续执行下一轮评估。

端到端验证结果表明，
当前原型能够完整执行：

\begin{equation}
\text{意图理解}
\rightarrow
\text{网络目标量化}
\rightarrow
\text{候选策略生成}
\rightarrow
\text{策略评估}
\rightarrow
\text{反馈改进}.
\label{eq:validated_closed_loop}
\end{equation}

这说明本文构建的四个模块之间的数据接口、
状态控制和异常处理能够形成可运行的闭环流程。


\subsection*{不同负载条件下的策略适应性}

为分析网络负载变化对候选策略可行性和性能的影响，
本文设置低负载、中负载和高负载三类RAN状态。

不同负载条件下的主要实验结果如表
\ref{tab:load_sweep_results}
所示。

\begin{table}[H]
    \centering
    \caption{不同网络负载条件下的策略评估结果}
    \label{tab:load_sweep_results}
    \small
    \renewcommand{\arraystretch}{1.25}
    \begin{tabular}{
        >{\centering\arraybackslash}p{0.11\textwidth}
        >{\centering\arraybackslash}p{0.14\textwidth}
        >{\centering\arraybackslash}p{0.14\textwidth}
        >{\centering\arraybackslash}p{0.17\textwidth}
        >{\centering\arraybackslash}p{0.16\textwidth}
        >{\centering\arraybackslash}p{0.13\textwidth}
    }
        \toprule
        负载状态
        & 可行率
        & 最优评分
        & P05吞吐率/Mbps
        & P95时延/ms
        & PRB利用率
        \\
        \midrule

        低负载
        & 0.95
        & 0.765702
        & 33.750
        & 15.357
        & 0.4375
        \\

        中负载
        & 0.87
        & 0.687047
        & 27.000
        & 21.500
        & 0.6575
        \\

        高负载
        & 0.14
        & 0.613658
        & 22.275
        & 29.179
        & 0.8375
        \\

        \bottomrule
    \end{tabular}
\end{table}

随着负载从低负载升高至高负载，
可行策略比例由0.95下降至0.14，
说明当前离散动作空间内能够同时满足视频、
语音和PRB硬约束的策略数量明显减少。

可行率下降可以表示为：

\begin{equation}
\Delta R_{\mathrm{feasible}}
=
0.14-0.87
=
-0.73,
\label{eq:feasible_rate_drop}
\end{equation}

即从中负载到高负载，
可行率下降0.73。

同时，
视频P05吞吐率由33.750 Mbps下降至22.275 Mbps，
视频P95时延由15.357 ms上升至29.179 ms，
PRB利用率由0.4375上升至0.8375。

这些变化与性能模型的预期方向一致：

\begin{equation}
u_0\uparrow
\quad\Rightarrow\quad
u'\uparrow,
\quad
o\uparrow,
\quad
T_{P05}\downarrow,
\quad
D_{P95}\uparrow.
\label{eq:load_performance_trend}
\end{equation}

不同负载下的Pareto候选数量分别为：

\begin{equation}
N_{\mathrm{Pareto}}
=
\left\{
11,9,4
\right\},
\label{eq:pareto_size_under_load}
\end{equation}

对应低、中、高三类负载。

这表明随着网络负载升高，
不仅满足全部硬约束的策略数量减少，
业务、能耗、资源和风险之间可供选择的非支配策略空间也明显收缩。

需要注意的是，
三类负载条件下的最优策略标识均为：

\begin{equation}
P^{*}
=
\texttt{RAN-POLICY-R00-0016}.
\label{eq:stable_best_policy}
\end{equation}

这说明在当前离散动作空间和评分模型下，
同一策略在三个预定义负载状态中均保持首名。

但是，这一结果只能说明当前候选集内的排序稳定性，
不能证明该策略在连续动作空间或真实网络中具有全局最优性。


\subsection*{反馈机制对比}

为分析反馈优化的实际作用，
本文分别在相同输入条件下运行：

\begin{equation}
feedbackEnabled
=
\texttt{false}
\end{equation}

和：

\begin{equation}
feedbackEnabled
=
\texttt{true}.
\end{equation}

反馈关闭时，
系统只执行第0轮候选生成和评估；
反馈开启时，
系统根据评估结果生成局部邻域策略并继续迭代。

对比结果如表
\ref{tab:feedback_comparison}
所示。

\begin{table}[H]
    \centering
    \caption{反馈机制开启与关闭的结果对比}
    \label{tab:feedback_comparison}
    \small
    \renewcommand{\arraystretch}{1.25}
    \begin{tabular}{
        >{\raggedright\arraybackslash}p{0.30\textwidth}
        >{\centering\arraybackslash}p{0.22\textwidth}
        >{\centering\arraybackslash}p{0.22\textwidth}
    }
        \toprule
        指标 & 反馈关闭 & 反馈开启 \\
        \midrule

        生成轮数
        & 1
        & 5
        \\

        反馈迭代次数
        & 0
        & 4
        \\

        评估策略数
        & 100
        & 109
        \\

        可行策略数
        & 14
        & 15
        \\

        可行率
        & 0.140000
        & 0.137615
        \\

        最优综合评分
        & 不变
        & 不变
        \\

        最优策略
        & 相同
        & 相同
        \\

        \bottomrule
    \end{tabular}
\end{table}

反馈开启后，
评估策略数量由100增加到109，
可行策略数量由14增加到15。

新发现的可行策略为：

\begin{equation}
P_{\mathrm{new}}
=
\texttt{RAN-POLICY-R02-FB-F73FCAE57E}.
\label{eq:new_feasible_feedback_policy}
\end{equation}

这说明当前反馈机制能够围绕父策略扩展局部搜索，
并发现初始模板网格中未保留的新可行策略。

但是，开启反馈后最优综合评分没有提升，
首名策略也没有发生变化。

因此，当前实验能够支持的结论是：

\begin{equation}
\boxed{
\text{反馈机制扩展了局部搜索空间并发现额外可行策略，
但未提高当前场景下的最优策略性能。}
}
\label{eq:feedback_experiment_conclusion}
\end{equation}

反馈开启后可行率略有下降，
并不是可行策略数量减少，
而是由于总评估策略数增加：

\begin{equation}
\frac{15}{109}
<
\frac{14}{100}.
\label{eq:feedback_feasible_ratio_explanation}
\end{equation}

因此，评价反馈机制时，
不能只比较可行率，
还需要同时比较可行策略绝对数量、最优评分和首名策略变化。

此外，
各实验种子下的结果完全一致，
原因是当前性能模型和候选生成过程均为确定性模式，
随机种子只用于结果追踪，
没有参与性能扰动。


\subsection*{评分权重敏感性}

为分析综合评分结果对权重配置的敏感程度，
本文设置以下五组权重方案：

\begin{itemize}
    \item 默认权重；
    \item 业务保障优先；
    \item 节能优先；
    \item 资源效率优先；
    \item 风险控制优先。
\end{itemize}

综合评分统一表示为：

\begin{equation}
Score(P_i)
=
\lambda_sS_{\mathrm{service}}
+
\lambda_eS_{\mathrm{energy}}
+
\lambda_rS_{\mathrm{resource}}
+
\lambda_qS_{\mathrm{risk}}.
\label{eq:weight_sensitivity_score}
\end{equation}

权重敏感性主要通过以下指标评价：

\begin{itemize}
    \item 最优策略是否变化；
    \item Top-5策略集合Jaccard相似度；
    \item 全部可行策略排名的Spearman相关系数；
    \item 相对于默认权重的排名变化数量。
\end{itemize}

主要结果如表
\ref{tab:weight_sensitivity_results}
所示。

\begin{table}[H]
    \centering
    \caption{评分权重敏感性结果}
    \label{tab:weight_sensitivity_results}
    \small
    \renewcommand{\arraystretch}{1.25}
    \begin{tabular}{
        >{\raggedright\arraybackslash}p{0.26\textwidth}
        >{\centering\arraybackslash}p{0.19\textwidth}
        >{\centering\arraybackslash}p{0.20\textwidth}
        >{\centering\arraybackslash}p{0.18\textwidth}
    }
        \toprule
        权重方案
        & Top-5 Jaccard
        & Spearman相关系数
        & 排名变化数
        \\
        \midrule

        默认权重
        & 1.0000
        & 1.0000
        & 0
        \\

        业务保障优先
        & 1.0000
        & 0.9824
        & 8
        \\

        节能优先
        & 1.0000
        & 1.0000
        & 0
        \\

        资源效率优先
        & 0.6667
        & 0.9604
        & 7
        \\

        风险控制优先
        & 1.0000
        & 1.0000
        & 0
        \\

        \bottomrule
    \end{tabular}
\end{table}

五组权重方案下的首名策略均为：

\begin{equation}
P^{*}
=
\texttt{RAN-POLICY-R00-0016}.
\label{eq:weight_stable_best_policy}
\end{equation}

说明在当前候选集和评分结构下，
最优策略对权重调整具有较高稳定性。

业务保障优先方案的Spearman相关系数为0.9824，
共出现8个策略排名变化，
说明提高业务评分权重会影响中上游策略的相对顺序，
但未改变Top-5策略集合。

资源效率优先方案的Top-5 Jaccard相似度为0.6667，
Spearman相关系数为0.9604，
共出现7个排名变化。
这表明资源权重变化对前列策略集合的影响相对更明显。

节能优先和风险控制优先方案下，
完整排序与默认权重保持一致。
这说明在当前性能结果和归一化边界下，
节能和风险评分差异不足以改变策略间的排序关系。

需要指出的是，
不同权重方案下的综合评分数值不能直接解释为物理性能提升或下降。
因为：

\begin{equation}
Score_{\lambda^{(1)}}(P)
\neq
Score_{\lambda^{(2)}}(P)
\label{eq:score_not_directly_comparable}
\end{equation}

时，
分值差异可能仅由评分偏好变化引起，
而不是策略实际吞吐率、时延或能耗发生变化。

实验还表明，
当前首名策略属于Pareto候选集合，
但并不支配全部可行策略。
因此，首名稳定主要来自当前评分结构、
候选动作空间和性能结果分布，
不能解释为该策略在所有评价维度上均优于其他策略。


\subsection*{模型裁剪与结果有效性边界}

为保证简化性能模型输出保持在物理或评分允许范围内，
系统对吞吐率、时延、BLER、语音质量、
能耗、普通用户下降和风险等指标执行边界裁剪：

\begin{equation}
x'
=
\operatorname{clip}
\left(
x_{\mathrm{raw}},
x_{\min},
x_{\max}
\right).
\label{eq:simulation_output_clipping}
\end{equation}

正式实验共记录4633次裁剪事件。

其中：

\begin{itemize}
    \item 风险归一化裁剪3644次，占78.65\%；
    \item 语音质量输出裁剪725次；
    \item 用户视频上行BLER裁剪264次。
\end{itemize}

裁剪事件数量关系为：

\begin{equation}
3644+725+264
=
4633.
\label{eq:clipping_event_count}
\end{equation}

风险归一化裁剪主要发生在风险原始值低于下界时，
表示部分候选策略的某些风险分量被归一化为0。

语音质量裁剪主要发生在模型原始计算结果高于1时，
系统将其限制为质量上界1。

BLER裁剪主要发生在资源保障和调度增益较强时，
原始BLER计算结果低于0，
系统将其限制为0。

为判断裁剪是否影响实验结论，
系统同时比较裁剪前后的：

\begin{itemize}
    \item 硬约束满足情况；
    \item 每轮首名策略；
    \item 最终推荐策略；
    \item 策略排序变化。
\end{itemize}

检查结果表明：

\begin{equation}
N_{\mathrm{constraint\ change}}
=
0,
\label{eq:no_constraint_change_by_clipping}
\end{equation}

即没有任何策略因裁剪导致硬约束判定发生变化。

同时：

\begin{equation}
N_{\mathrm{top\ change}}
=
0,
\label{eq:no_top_policy_change_by_clipping}
\end{equation}

即裁剪没有改变每轮首名策略和最终推荐策略。

部分非首名策略的相对顺序发生变化，
但不影响本章关于负载趋势、
反馈作用和首名稳定性的主要结论。

因此，当前裁剪现象属于需要记录的模型诊断信息，
但在本次实验中不构成阻断性问题。

需要强调的是，
裁剪比例较高也说明当前简化模型的部分公式和归一化边界
仍有进一步标定和优化空间。
未来接入高精度网络数字孪生后，
应根据真实网络数据重新校准模型参数和输出范围。


\subsection*{实验结果讨论}

综合上述实验结果，
可以得到以下主要结论。

第一，
当前原型能够稳定完成从自然语言业务需求到RAN策略推荐的完整闭环，
说明意图抽取、意图转译、策略生成和策略评估之间的数据接口设计有效。

第二，
随着网络负载升高，
视频P05吞吐率下降、P95时延上升、PRB利用率增加，
可行策略比例和Pareto候选数量均明显下降。
这表明高负载条件下，
同时满足视频保障、语音保护和资源约束的策略空间显著收缩。

第三，
反馈机制能够在当前最优或近似最优父策略周围
生成新的局部候选，
并发现额外可行策略。
但是，在当前场景和离散动作空间下，
反馈并未提升最优综合评分，
说明局部邻域搜索的改进空间有限。

第四，
评分权重变化没有改变首名策略，
但业务权重和资源权重变化会影响部分中上游策略排序。
因此，在实际系统中，
评分权重仍需要根据业务优先级和运营目标进行配置与校准。

第五，
模型裁剪没有改变硬约束判定和最终推荐结果，
但较高的风险归一化裁剪比例表明，
当前简化模型仍需要进一步标定。

上述结论适用于：

\begin{itemize}
    \item \texttt{simplified-ran-v1}性能模型；
    \item 当前单小区RAN状态；
    \item 当前四类RAN动作；
    \item 当前离散参数集合；
    \item 当前SLA门限和评分权重；
    \item 当前预定义的低、中、高负载状态。
\end{itemize}

不能将实验结果直接外推到：

\begin{itemize}
    \item 真实运营商网络；
    \item 多小区和跨域网络；
    \item 连续动作参数空间；
    \item 不同业务负载和用户分布；
    \item 真实无线信道、调度和队列过程；
    \item 真实基站能耗和设备控制。
\end{itemize}

因此，本节实验的定位是：

\begin{equation}
\boxed{
\text{验证意图驱动决策流程和相对变化趋势，
而不是预测真实RAN性能。}
}
\label{eq:simulation_result_positioning}
\end{equation}

当前性能评估器已经通过统一接口与上层策略决策模块解耦。
下一章可在保持意图转译、
策略生成和评分接口基本不变的前提下，
将\texttt{simplified-ran-v1}替换为
具有更高保真度的网络数字孪生评估器，
进一步提升策略推演结果的可信度。

\chapter{网络高精度孪生建模与编排}

网络高精度孪生建模与编排是支撑6G意图驱动智能决策可信验证的重要技术基础。第三章中的意图驱动智能决策主要解决“业务意图如何转换为候选网络策略”的问题，而本章重点解决“如何构建可信数字孪生环境，并利用孪生环境对候选策略进行仿真推演和效果验证”的问题。

在本文语境下，网络高精度孪生建模与编排并不侧重现网中的网络功能编排或资源下发，而是侧重数字孪生环境内部的数据、模型、场景、仿真任务和验证流程组织。也就是说，建模解决“数字孪生环境如何准确刻画真实网络”的问题，编排解决“如何将数据、模型和仿真任务组织起来支撑策略验证”的问题。

本章围绕网络高精度孪生建模与编排关键技术展开，主要包括多源数据融合、网络语义建模、高精度孪生模型构建、孪生模型动态校准、数字孪生环境编排以及基于孪生环境的策略仿真验证等内容。

\section{网络高精度孪生建模与编排概述}

网络数字孪生（Network Digital Twin, NDT）的核心思想是在数字空间中构建物理网络的虚拟映射体，并通过数据交互和模型更新实现对物理网络状态、行为和性能的动态刻画。面向6G网络，网络数字孪生不仅需要反映网络当前状态，还需要具备性能预测、策略推演和风险评估能力，为智能决策提供可信验证环境。

网络高精度孪生建模与编排可分为两个层面。

一是高精度孪生建模，主要面向物理网络状态和行为的数字化表达，重点构建网络拓扑、无线环境、业务流量、资源状态、协议流程和性能预测等模型，使数字孪生环境能够尽可能准确地反映真实网络。

二是数字孪生环境编排，主要面向特定业务意图和候选策略验证任务，对孪生环境中的数据资源、模型资源、场景资源、计算资源和仿真任务进行组织与调度，形成可运行、可验证、可评估的数字孪生实验环境。

二者关系可以概括为：

\begin{equation}
Physical\ Network
\rightarrow
Twin\ Modeling
\rightarrow
Twin\ Orchestration
\rightarrow
Simulation
\rightarrow
Evaluation
\end{equation}

其中，$Twin\ Modeling$表示网络孪生建模过程，$Twin\ Orchestration$表示数字孪生环境中的模型、场景和任务编排过程，$Simulation$表示仿真推演过程，$Evaluation$表示策略效果评估过程。

\section{多源异构数据感知与融合}

高精度孪生建模首先依赖多源异构数据。6G网络涉及无线接入、核心网、承载网、边缘计算、终端和业务平台等多个对象，其数据来源复杂、类型多样、时间粒度不一致，需要通过统一的数据治理和融合机制，为孪生建模提供可靠数据基础。

面向网络数字孪生的数据可表示为：

\begin{equation}
D=\{D_{topo},D_{config},D_{KPI},D_{service},D_{resource},D_{env}\}
\end{equation}

其中，$D_{topo}$表示网络拓扑数据，$D_{config}$表示配置参数数据，$D_{KPI}$表示性能指标数据，$D_{service}$表示业务数据，$D_{resource}$表示资源状态数据，$D_{env}$表示环境数据。

具体包括：

\begin{itemize}
    \item \textbf{网络拓扑数据：} 描述基站、小区、核心网网元、边缘节点、链路和接口关系；
    \item \textbf{配置参数数据：} 描述小区频点、带宽、功率、邻区关系、切片参数等配置状态；
    \item \textbf{性能指标数据：} 包括吞吐量、时延、丢包率、资源利用率、接入成功率、切换成功率等；
    \item \textbf{业务状态数据：} 描述业务类型、业务流量、用户分布、业务体验和SLA状态；
    \item \textbf{资源状态数据：} 包括频谱、功率、PRB、算力、存储和传输资源等；
    \item \textbf{环境数据：} 描述地理位置、无线传播环境、干扰状态、用户移动轨迹等。
\end{itemize}

通过数据清洗、时空对齐、特征提取和关联融合，可形成面向孪生建模的统一数据视图，为后续语义建模、状态预测和策略验证提供基础。

\section{基于知识图谱的网络语义建模}

仅依靠结构化指标难以完整表达网络实体之间复杂的关联关系。为了提升孪生模型的可解释性和可扩展性，需要引入知识图谱对网络实体、业务对象、性能指标、资源状态和控制策略之间的语义关系进行建模。

网络知识图谱可表示为：

\begin{equation}
G=(V,E,A)
\end{equation}

其中，$V$表示实体集合，$E$表示实体之间的关系集合，$A$表示实体属性集合。

在6G网络数字孪生场景中，典型实体包括：

\begin{itemize}
    \item 网络实体：基站、小区、UE、核心网网元、边缘节点等；
    \item 资源实体：频谱、功率、PRB、算力、存储、链路等；
    \item 业务实体：工业控制、高清视频、车联网、低空通信等；
    \item 指标实体：时延、吞吐量、可靠性、丢包率、能耗、QoE等；
    \item 策略实体：资源调度、功率调整、切片保障、节能控制等。
\end{itemize}

典型关系包括：

\begin{itemize}
    \item 连接关系：如“UE连接小区”“基站连接核心网网元”；
    \item 影响关系：如“干扰增强导致误块率升高”；
    \item 约束关系：如“业务SLA约束资源分配策略”；
    \item 因果关系：如“PRB不足导致吞吐量下降”；
    \item 映射关系：如“业务意图映射为SLA指标”。
\end{itemize}

通过知识图谱，可实现从网络数据到网络知识的转换，为孪生模型构建、策略验证和结果解释提供语义支撑。

\section{高精度孪生模型构建}

高精度孪生模型是数字孪生环境的核心。面向6G网络，单一模型难以完整描述网络状态，需要采用多模型融合方式，构建覆盖拓扑、无线、业务、资源、协议和AI预测能力的模型体系。

\subsection{拓扑模型}

拓扑模型用于描述网络节点、链路和连接关系，可表示为：

\begin{equation}
T=(N,L)
\end{equation}

其中，$N$表示网络节点集合，$L$表示链路集合。拓扑模型为业务路径分析、资源关联分析和故障影响范围评估提供基础。

\subsection{无线环境模型}

无线环境模型用于描述信道状态、覆盖范围、干扰关系和用户移动行为。其可抽象表示为：

\begin{equation}
H_t=f(d_t,\theta_t,v_t,\rho_t)
\end{equation}

其中，$d_t$表示传播距离，$\theta_t$表示空间方向，$v_t$表示用户移动状态，$\rho_t$表示环境参数。无线环境模型是无线接入网络孪生建模中的关键模型。

\subsection{业务行为模型}

业务行为模型用于描述业务流量到达、用户行为和业务体验变化，可表示为：

\begin{equation}
B_t=f(\lambda_t,U_t,QoS_t)
\end{equation}

其中，$\lambda_t$表示业务到达率，$U_t$表示用户状态，$QoS_t$表示业务质量需求。

\subsection{资源状态模型}

资源状态模型用于描述网络资源占用和变化情况，包括无线资源、计算资源、传输资源和存储资源等，可表示为：

\begin{equation}
R_t=\{R_{radio},R_{compute},R_{transport},R_{storage}\}
\end{equation}

该模型为资源调度、策略验证和资源冲突分析提供依据。

\subsection{AI预测模型}

AI预测模型用于根据历史状态和实时数据预测未来网络状态变化：

\begin{equation}
\hat{X}_{t+k}=g_{\phi}(X_{0:t},D_{0:t})
\end{equation}

其中，$\hat{X}_{t+k}$表示未来时刻的预测网络状态，$g_{\phi}$表示AI预测模型，$X_{0:t}$和$D_{0:t}$分别表示历史网络状态和历史观测数据。

通过上述多模型融合，可以形成面向策略验证的高精度网络孪生模型体系。

\section{孪生模型动态校准与一致性维护}

由于真实网络状态具有动态变化特征，数字孪生模型需要持续根据物理网络反馈进行动态校准，保证虚拟网络与物理网络之间的一致性。

孪生误差可表示为：

\begin{equation}
e_t=\|X_t-\hat{X}_t\|
\end{equation}

其中，$X_t$表示物理网络真实状态，$\hat{X}_t$表示数字孪生模型估计状态，$e_t$表示孪生误差。

模型校准目标是最小化孪生误差：

\begin{equation}
\min_{\theta}\sum_{t=1}^{T}\|X_t-f_{\theta}(D_{0:t},M_t)\|^2
\end{equation}

其中，$f_{\theta}$表示孪生状态映射函数，$D_{0:t}$表示历史观测数据，$M_t$表示模型集合。

动态校准主要包括：

\begin{itemize}
    \item 状态同步：根据实时网络数据更新孪生状态；
    \item 参数修正：根据观测误差调整模型参数；
    \item 模型更新：根据网络结构和业务变化更新模型；
    \item 版本管理：记录模型演进过程，支持回溯与比较；
    \item 一致性评估：评估数字孪生环境与物理网络之间的偏差。
\end{itemize}

通过动态校准与一致性维护，数字孪生模型能够持续逼近物理网络状态，提高策略验证结果的可信度。

\section{数字孪生环境编排}

数字孪生环境编排是本文中“编排”的重点。其核心不是将策略直接下发到现网，而是根据业务意图和策略验证任务，对数字孪生环境中的数据、模型、场景、计算资源和仿真流程进行组织与调度，形成可运行的孪生验证环境。

数字孪生环境编排可表示为：

\begin{equation}
O_{twin}=F(D,M,Sce,Task,Res)
\end{equation}

其中，$D$表示数据资源，$M$表示模型资源，$Sce$表示仿真场景，$Task$表示仿真任务，$Res$表示计算资源。

具体包括以下内容：

\subsection{数据编排}

数据编排负责根据验证任务选择所需数据，例如拓扑数据、业务数据、性能数据、资源数据和环境数据，并完成数据加载、时间对齐和格式转换。

\subsection{模型编排}

模型编排负责根据场景需求选择和组合相应孪生模型。例如，工业控制低时延场景需要组合无线环境模型、业务时延模型、可靠性模型和资源调度模型；节能场景则需要组合流量预测模型、能耗模型和覆盖保障模型。

\subsection{场景编排}

场景编排负责构建面向特定业务目标的孪生实验场景，包括区域拓扑、用户分布、业务类型、资源状态、SLA约束和候选策略配置。

\subsection{任务编排}

任务编排负责组织仿真任务执行过程，包括任务创建、模型加载、参数配置、并行仿真、结果汇聚和指标评估等。

\subsection{资源编排}

资源编排负责调度数字孪生环境中的计算资源和仿真资源，保障仿真任务能够按照优先级和时延要求完成。

通过上述编排过程，数字孪生环境能够根据不同业务意图和策略验证需求，快速构建对应的仿真验证空间。

\section{基于孪生环境的策略仿真验证}

第三章输出的候选策略需要在数字孪生环境中进行仿真验证。该过程通过在孪生空间中模拟候选策略执行效果，评估其是否满足SLA要求，并分析潜在风险。

策略验证流程如下：

\begin{enumerate}
    \item 输入候选策略集合；
    \item 调用数字孪生环境编排机制，构建验证场景；
    \item 加载相关数据、模型和SLA约束；
    \item 在孪生环境中执行仿真推演；
    \item 评估策略对性能、体验、资源和能耗的影响；
    \item 输出策略验证结果和优化建议。
\end{enumerate}

策略评价函数可表示为：

\begin{equation}
J(P_i)=\alpha QoE(P_i)-\beta Delay(P_i)+\gamma Reliability(P_i)-\eta Energy(P_i)
\end{equation}

其中，$P_i$表示候选策略，$QoE(P_i)$表示策略对应的用户体验，$Delay(P_i)$表示时延，$Reliability(P_i)$表示可靠性，$Energy(P_i)$表示能耗，$\alpha$、$\beta$、$\gamma$、$\eta$为权重系数。

最终，数字孪生环境输出策略验证结果：

\begin{equation}
Result=\{P_i,J(P_i),Risk(P_i),SLA(P_i)\}
\end{equation}

其中，$Risk(P_i)$表示策略风险，$SLA(P_i)$表示策略对SLA约束的满足情况。

通过该过程，数字孪生环境能够在策略作用于真实网络之前完成预验证，降低现网试错风险，提高智能决策的可信性。

\section{小结}

本章围绕网络高精度孪生建模与编排技术展开，重点研究了多源数据融合、网络语义建模、高精度孪生模型构建、模型动态校准、数字孪生环境编排和策略仿真验证等内容。

其中，高精度孪生建模侧重于构建物理网络在数字空间中的准确表达，解决“孪生体如何准确刻画真实网络”的问题；数字孪生环境编排侧重于根据业务意图和策略验证任务，对数据、模型、场景、计算资源和仿真任务进行组织调度，解决“孪生体如何被有效调用并支撑策略验证”的问题。

整体来看，本章解决的是“策略是否有效、是否安全、如何验证”的问题，与第三章“意图驱动智能决策技术”形成衔接关系，共同支撑面向6G自智网络的智能决策与可信验证闭环。




\printbibliography
\end{document}
